\section{The general model and its fundamental equation}\label{sec:model_x}

This section defines the general two-class system with both jockeying and abandonment, which we call \emph{Model-$X$}; the simpler variants studied in the later sections are recovered from it as special cases.

The model is an $M/M/1$ queue with two priority classes. Customers are served First-Come, First-Served (FCFS) within their own class, and class-$1$ (the priority class) has non-preemptive priority over class-$2$: a service in progress is never interrupted, so a class-$2$ customer in service is not displaced by an arriving class-$1$ customer. The server is work-conserving: it is never idle while a customer is present. Beyond these basic dynamics, Model-$X$ adds two queue-leaving mechanisms. First, \emph{jockeying}: a waiting customer leaves its own queue to \emph{join the other queue}, at per-customer rate $\gamma_1$ for direction $1\to2$ and $\gamma_2$ for $2\to1$, so the total number of customers is preserved. Second, \emph{abandonment} (reneging): a waiting customer leaves the \emph{system} at per-customer rate $\theta_i$, decreasing the total by one. Concretely, each waiting class-$i$ customer (not the one in service) holds an independent $\mathrm{Exp}(\theta_i)$ patience clock and abandons when it expires. This is precisely the linear-reneging mechanism of the $M/M/1+M$ queue from \S\ref{ssec:erlangA}.

In Kendall's notation this is an $M/M/1$-type model: Poisson arrivals, exponentially distributed service times, and a single server. Throughout, the parameters are:
\begin{itemize}
    \item $\lambda_1, \lambda_2$: Poisson arrival rates of classes $1$ and $2$, respectively; the total arrival rate is $\lambda=\lambda_1+\lambda_2$;
    \item $\mu$: exponential service rate of a customer, regardless of class;
    \item $\rho_i = \frac{\lambda_i}{\mu}$: system load of class $i \in \{1,2\}$; the total system load is $\rho=\frac{\lambda}{\mu}$;
    \item $\gamma_1, \gamma_2$: per-customer jockeying rates from Class $1\to2$ and $2\to1$, respectively; jockeying transfers a customer between queues, preserving the total number of customers;
    \item $\theta_1, \theta_2$: per-customer abandonment (reneging) rates from classes $1$ and $2$, respectively; abandonment removes a customer from the system, reducing the total by one.
\end{itemize}
The two class arrival streams are independent Poisson processes, and they are independent of the system's internal exponential service, jockeying and abandonment clocks, so they satisfy the lack-of-anticipation assumption of the PASTA property (\S\ref{ssec:priority}). By Poisson \emph{superposition} the merged arrival stream is itself Poisson at rate $\lambda$, and each class stream is an i.i.d.\ $\lambda_i/\lambda$ \emph{thinning} of the merger, so PASTA applies to the aggregate state and to each class separately.

To define the stochastic processes that drive this system, we introduce the following random variables:
\begin{itemize}
    \item $\mathcal{B}:$ binary indicator for the server being idle or busy, $\mathcal{B}\in \{0,1\}$;
    \item $N_1, N_2:$ number of customers of class $1$ and $2$, respectively, in their respective queues;
    \item $N:$ total number of customers in the queues.
\end{itemize}
Note that $N_1$, $N_2$ and $N$ count customers in the \emph{queues}, not the \emph{system}. Because the service rate $\mu$ is common to both classes, the \emph{memorylessness} of the exponential service time makes the residual service of the customer in progress $\mathrm{Exp}(\mu)$ regardless of its class; the class of the in-service customer therefore does not influence the future evolution of the queue lengths, and the state descriptor need not record it. Accordingly, we write the idle state as $(0)$ and all other states as $(n_1,n_2)$, so that $(0,0)$ implies that the server is busy while both queues are empty. Formally,
\begin{equation}
    S := \{(0)\} \cup \{ (n_1,n_2) \mid n_1,n_2 \geq 0 \}.
    \label{eq:gen:state_space_S}
\end{equation}
For the partial-PGF analysis we additionally use an auxiliary state space $\widetilde{S}$, which tracks only $n_2$ and the total $n = n_1 + n_2$, subject to $n \geq n_2 \geq 0$:
\begin{equation}
    \widetilde{S} := \{(0)\} \cup \{ (n_2,n) \mid n \geq n_2 \geq 0 \}.
    \label{eq:gen:state_space_Stilde}
\end{equation}

Assuming positive recurrence, so that the limiting distribution exists, we define the long-run probabilities of finding the system in each state:
\begin{itemize}
    \item $\pi_0$ for the idle state $(0)$, it belongs to both spaces $S$ and $\widetilde{S}$;
    \item $\pi(n_1,n_2)$ for states $(n_1,n_2) \in S$, where $n_1,n_2 \geq 0$;
    \item $\widetilde{\pi}(n_2,n)$ for states $(n_2,n) \in \widetilde{S}$, where $n \geq n_2 \geq 0$.
\end{itemize}

The simpler variants are recovered from Model-$X$ by setting the appropriate parameters to zero, as summarised in Table~\ref{tab:model_family}:
\begin{table}[H]
    \centering
    \begin{tabular}{cccl}
        \hline
        Model & Jockeying & Abandonment & Parameter values \\
        \hline
        $X$  & both ($1\!\leftrightarrow\!2$) & both ($1,2$)       & $\gamma_i \geq 0,\ \theta_i \geq 0$ \quad (general) \\
        $A$  & none                            & none               & $\gamma_1=\gamma_2=0,\ \theta_1=\theta_2=0$ \\
        $B$  & both ($1\!\leftrightarrow\!2$) & none               & $\gamma_i \geq 0,\ \theta_1=\theta_2=0$ \\
        $B_2$ & $1\!\to\!2$ only                & none               & $\gamma_1 \geq 0,\ \gamma_2=0,\ \theta_1=\theta_2=0$ \\
        $C_2$ & none                            & class $1$ only     & $\theta_1 \geq 0,\ \theta_2=0,\ \gamma_1=\gamma_2=0$ \\
        \hline
        $\BH$ & $1\!\to\!2$, head-of-line       & none                    & $\gamma_1\,\mathbf{1}_{\{n_1\ge1\}}>0,\ \gamma_2=\theta_1=\theta_2=0$ \\
        $\CH$ & none                            & class $1$, head-of-line & $\theta_1\,\mathbf{1}_{\{n_1\ge1\}}>0,\ \gamma_1=\gamma_2=\theta_2=0$ \\
        \hline
    \end{tabular}
    \caption{Nested family of models. Model-$X$ is the parent; Models~$A$, $B$, $B_2$ and $C_2$
    are obtained by restricting the jockeying directions and the abandoning classes. The two
    head-of-line variants $\BH$ and $\CH$ are listed separately, as their mechanism involves the indicator $\mathbf{1}_{\{n_1\ge1\}}$ rather than being proportional to the queue length.}
    \label{tab:model_family}
    \par\smallskip
\end{table}

Figure~\ref{fig:models-queue-overview} gives a visual overview of Model-$X$ and each variant we consider. It omits the two head-of-line variants, which are less conventional and for which such a diagram would add little:
\begin{figure}[H]
    \centering
    \resizebox{\textwidth}{!}{%
        % Two-column overview of all five models.
%
% Left column  : Model X (scale 0.75), vertically spanning the full grid height.
% Right column : 2×2 grid of submodels (scale 0.38), row-pitch 3.2 units.
%
%   [ Model X ]  [ Model A  |  Model B  ]
%   [         ]  [ Model B₂ |  Model C₂ ]
%
% No legend — arrow conventions are explained in the caption.
%
\begin{tikzpicture}[>=Latex]

% ─────────────────────────────────────────────────────────────────────────────
% LEFT COLUMN: Model X
%
% Local coords: x ∈ [−1.5, 7], y ∈ [−3, 3]  (span 8.5 × 6)
% Scale 0.75 → physical 6.375 × 4.5.
% Vertical centre of right grid (see below): y_centre ≈ 1.6
% → shift_y = 1.6  so local y=0 (server) aligns with grid midpoint.
% → shift_x = 1.0  so left edge of diagram lands near x=0.
% ─────────────────────────────────────────────────────────────────────────────
\begin{scope}[shift={(1.0, 1.6)}, scale=0.75, line width=0.7pt]
    % Queue 1 (top, Class-1, priority)
    \draw (0,1.9)--(2,1.9)--(2,1.1)--(0,1.1);
    \foreach \x in {0.2,0.4,...,1.8}{\draw[line width=0.35pt](\x,1.8)--(\x,1.2);}
    \draw[->](-1.5,1.5)--(-0.1,1.5) node[midway,above,font=\small]{$\lambda_1$};
    \draw[->,dotted](1,2.0)--(1,3.0) node[above,font=\small]{$\theta_1$};
    % Queue 2 (bottom, Class-2)
    \draw (0,-1.1)--(2,-1.1)--(2,-1.9)--(0,-1.9);
    \foreach \x in {0.2,0.4,...,1.8}{\draw[line width=0.35pt](\x,-1.2)--(\x,-1.8);}
    \draw[->](-1.5,-1.5)--(-0.1,-1.5) node[midway,above,font=\small]{$\lambda_2$};
    \draw[->,dotted](1,-2.0)--(1,-3.0) node[below,font=\small]{$\theta_2$};
    % Jockeying
    \draw[->,dashed](0.6,1.0)--(0.6,-1.0) node[midway,left,font=\small]{$\gamma_1$};
    \draw[->,dashed](1.4,-1.0)--(1.4,1.0) node[midway,right,font=\small]{$\gamma_2$};
    % Server
    \node[draw,circle,minimum size=1.1cm,font=\normalsize](sX) at (5,0){$\mu$};
    \draw[->](2.2,1.5)--(sX);
    \draw[->](2.2,-1.5)--(sX);
    \draw[->](sX)--(7,0);
\end{scope}
% Label above Model X
\node[font=\normalsize\bfseries] at (3.1, 4.3) {Model~X (parent)};

% Thin vertical separator between columns
\draw[gray!50, thin] (6.8, -1.1) -- (6.8, 4.3);

% ─────────────────────────────────────────────────────────────────────────────
% RIGHT COLUMN: 2×2 grid of submodels (scale 0.38, line width 0.55pt)
%
% Row 1 (top):    A at x-shift 7.05,  B  at x-shift 10.85
% Row 2 (bottom): B₂ at x-shift 7.05, C₂ at x-shift 10.85
% Row y-shifts:   top row 3.2,  bottom row 0.0
% Vertical centre of grid: (−1.9×0.38 + 3.2 + 1.9×0.38) / 2 ≈ 1.6  ← matches Model X
% ─────────────────────────────────────────────────────────────────────────────

% ── Model A (top-left) ───────────────────────────────────────────────────────
\begin{scope}[shift={(7.05, 3.2)}, scale=0.38, line width=0.55pt]
    \draw (0,1.9)--(2,1.9)--(2,1.1)--(0,1.1);
    \foreach \x in {0.2,0.4,...,1.8}{\draw[line width=0.3pt](\x,1.8)--(\x,1.2);}
    \draw[->](-1.5,1.5)--(-0.1,1.5) node[midway,above,font=\tiny]{$\lambda_1$};
    \draw (0,-1.1)--(2,-1.1)--(2,-1.9)--(0,-1.9);
    \foreach \x in {0.2,0.4,...,1.8}{\draw[line width=0.3pt](\x,-1.2)--(\x,-1.8);}
    \draw[->](-1.5,-1.5)--(-0.1,-1.5) node[midway,above,font=\tiny]{$\lambda_2$};
    \node[draw,circle,minimum size=0.85cm,font=\small](sA) at (5,0){$\mu$};
    \draw[->](2.2,1.5)--(sA); \draw[->](2.2,-1.5)--(sA); \draw[->](sA)--(7,0);
\end{scope}
\node[font=\small\bfseries] at (8.10, 2.10) {Model~A};

% ── Model B (top-right) ──────────────────────────────────────────────────────
\begin{scope}[shift={(10.85, 3.2)}, scale=0.38, line width=0.55pt]
    \draw (0,1.9)--(2,1.9)--(2,1.1)--(0,1.1);
    \foreach \x in {0.2,0.4,...,1.8}{\draw[line width=0.3pt](\x,1.8)--(\x,1.2);}
    \draw[->](-1.5,1.5)--(-0.1,1.5) node[midway,above,font=\tiny]{$\lambda_1$};
    \draw (0,-1.1)--(2,-1.1)--(2,-1.9)--(0,-1.9);
    \foreach \x in {0.2,0.4,...,1.8}{\draw[line width=0.3pt](\x,-1.2)--(\x,-1.8);}
    \draw[->](-1.5,-1.5)--(-0.1,-1.5) node[midway,above,font=\tiny]{$\lambda_2$};
    \draw[->,dashed](0.6,1.0)--(0.6,-1.0) node[midway,left,font=\tiny]{$\gamma_1$};
    \draw[->,dashed](1.4,-1.0)--(1.4,1.0) node[midway,right,font=\tiny]{$\gamma_2$};
    \node[draw,circle,minimum size=0.85cm,font=\small](sB) at (5,0){$\mu$};
    \draw[->](2.2,1.5)--(sB); \draw[->](2.2,-1.5)--(sB); \draw[->](sB)--(7,0);
\end{scope}
\node[font=\small\bfseries] at (11.90, 2.10) {Model~B};

% ── Model B₂ (bottom-left) ───────────────────────────────────────────────────
\begin{scope}[shift={(7.05, 0.0)}, scale=0.38, line width=0.55pt]
    \draw (0,1.9)--(2,1.9)--(2,1.1)--(0,1.1);
    \foreach \x in {0.2,0.4,...,1.8}{\draw[line width=0.3pt](\x,1.8)--(\x,1.2);}
    \draw[->](-1.5,1.5)--(-0.1,1.5) node[midway,above,font=\tiny]{$\lambda_1$};
    \draw (0,-1.1)--(2,-1.1)--(2,-1.9)--(0,-1.9);
    \foreach \x in {0.2,0.4,...,1.8}{\draw[line width=0.3pt](\x,-1.2)--(\x,-1.8);}
    \draw[->](-1.5,-1.5)--(-0.1,-1.5) node[midway,above,font=\tiny]{$\lambda_2$};
    \draw[->,dashed](1.0,1.0)--(1.0,-1.0) node[midway,left,font=\tiny]{$\gamma_1$};
    \node[draw,circle,minimum size=0.85cm,font=\small](sB2) at (5,0){$\mu$};
    \draw[->](2.2,1.5)--(sB2); \draw[->](2.2,-1.5)--(sB2); \draw[->](sB2)--(7,0);
\end{scope}
\node[font=\small\bfseries] at (8.10, -1.10) {Model~$B_2$};

% ── Model C₂ (bottom-right) ──────────────────────────────────────────────────
\begin{scope}[shift={(10.85, 0.0)}, scale=0.38, line width=0.55pt]
    \draw (0,1.9)--(2,1.9)--(2,1.1)--(0,1.1);
    \foreach \x in {0.2,0.4,...,1.8}{\draw[line width=0.3pt](\x,1.8)--(\x,1.2);}
    \draw[->](-1.5,1.5)--(-0.1,1.5) node[midway,above,font=\tiny]{$\lambda_1$};
    \draw[->,dotted](1,2.0)--(1,3.0) node[above,font=\tiny]{$\theta_1$};
    \draw (0,-1.1)--(2,-1.1)--(2,-1.9)--(0,-1.9);
    \foreach \x in {0.2,0.4,...,1.8}{\draw[line width=0.3pt](\x,-1.2)--(\x,-1.8);}
    \draw[->](-1.5,-1.5)--(-0.1,-1.5) node[midway,above,font=\tiny]{$\lambda_2$};
    \node[draw,circle,minimum size=0.85cm,font=\small](sC2) at (5,0){$\mu$};
    \draw[->](2.2,1.5)--(sC2); \draw[->](2.2,-1.5)--(sC2); \draw[->](sC2)--(7,0);
\end{scope}
\node[font=\small\bfseries] at (11.90, -1.10) {Model~C$_2$};

\end{tikzpicture}
%
    }
    \caption{Physical queue layout for Model-$X$ (left) and its four specialisations (right, $2\times2$ grid). In all models, Queue~$1$ (top buffer) has non-preemptive priority over Queue~$2$. Solid arrows are for Poisson arrivals ($\lambda_i$) and exponential service ($\mu$); dashed arrows are for jockeying($\gamma_i$); and dotted arrows are for abandonments ($\theta_i$)
    }
    \label{fig:models-queue-overview}
\end{figure}

\noindent\textbf{Stability.} Jockeying does not affect stability (customers are merely relocated). Model-$X$ has abandonments from both classes, so the departure rate $\mu + \theta_1 n_1 + \theta_2 n_2$ grows unbounded, ensuring positive recurrence for all loads. In contrast, models without abandonments require $\rho < 1$ (they behave like M/M/1). Mixed cases, such as Model-$C_2$, are discussed separately (\S\ref{sec:model_c2}).

\medskip\noindent\textbf{Computation of $\pi_0$ and $\pi(0,0)$.} The two empty-state probabilities are determined, for the full Model-$X$, by two exact identities that hold irrespective of the jockeying and abandonment mechanisms: a cut between the idle and busy states, and a renewal--reward identity over busy cycles. Together they reduce the determination of $\pi_0$ and $\pi(0,0)$ to a single scalar, the mean busy period $\mathbb{E}[BP]$.

\begin{theorem}[Empty-state probabilities of the general model]\label{thm:gen:boundary}
    Assume the CTMC of Model-$X$ is positive recurrent, so that the stationary distribution exists. Let $\mathbb{E}[BP]$ denote the mean length of a server busy period ---the time from an arrival that finds the system empty until the system is next empty. Then
    \begin{equation}
        \pi(0,0) \;=\; \rho\,\pi_0,
        \label{eq:gen:pi00_general}
    \end{equation}
    and
    \begin{equation}
        \pi_0 \;=\; \frac{1}{1+\lambda\,\mathbb{E}[BP]},
        \qquad
        \pi(0,0) \;=\; \frac{\rho}{1+\lambda\,\mathbb{E}[BP]}.
        \label{eq:gen:pi0_renewal}
    \end{equation}
\end{theorem}

\begin{proof}
    \emph{Cut identity.} Consider the cut separating the idle state $(0)$ from the set of busy states. The only transitions out of $(0)$ are the class-$1$ and class-$2$ arrivals, at total rate $\lambda$; the only transition into $(0)$ is a service completion from $(0,0)$, at rate $\mu$ ---from any state with a waiting customer a departure leaves the server busy, and neither jockeying nor abandonment can act when both queues are empty. Equating the probability flux across the cut gives $\lambda\,\pi_0=\mu\,\pi(0,0)$, which is Equation~\eqref{eq:gen:pi00_general}.

    \emph{Renewal--reward identity.} The idle state is a regeneration point: by the memorylessness of the exponential inter-arrival and service times (the strong Markov property), the evolution after each visit to $(0)$ is an independent probabilistic copy of the process. Successive cycles therefore alternate an idle period of mean $1/\lambda$ (the wait for the next arrival, exponential at rate $\lambda$) and a busy period of mean $\mathbb{E}[BP]$. By the renewal--reward theorem, the stationary probability of the idle state equals the long-run fraction of time spent idle,
    \begin{equation*}
        \pi_0 \;=\; \frac{1/\lambda}{1/\lambda+\mathbb{E}[BP]} \;=\; \frac{1}{1+\lambda\,\mathbb{E}[BP]},
    \end{equation*}
    and substituting into Equation~\eqref{eq:gen:pi00_general} yields the stated $\pi(0,0)$.
\end{proof}

Both identities are exact for the full model. They reduce the problem to the mean busy period $\mathbb{E}[BP]$, which is elementary when there is no abandonment and is computed model by model once abandonment is present; the resulting $\pi_0$ and $\pi(0,0)$ for each tractable model then follow as corollaries (Corollary~\ref{cor:gen:no_abandon} below for the no-abandonment family; \S\ref{sec:model_c2} and \S\ref{sec:model_c21} for the abandonment variants). For the full asymmetric Model-$X$, $\mathbb{E}[BP]$ itself has no closed form (Remark~\ref{rem:B:modelX}).

The no-abandonment case is concrete enough to state at once; it covers Models~$A$, $B$, $B_2$ and $\BH$ simultaneously.

\begin{corollary}[Empty states without abandonment]\label{cor:gen:no_abandon}
    Suppose $\theta_1=\theta_2=0$ (Models~$A$, $B$, $B_2$ and $\BH$) and $\rho<1$. Then the total number in the system is a birth--death process with birth rate $\lambda$ and death rate $\mu$ ---an $M/M/1$ queue--- for \emph{any} jockeying rates $\gamma_1,\gamma_2$, since jockeying merely relocates a waiting customer and so conserves the total count. Its mean busy period is $\mathbb{E}[BP]=(\mu-\lambda)^{-1}$, and Theorem~\ref{thm:gen:boundary} gives
    \begin{equation}
        \pi_0 = 1-\rho, \qquad \pi(0,0)=\rho(1-\rho).
        \label{eq:gen:pi0_noabandon}
    \end{equation}
\end{corollary}

\subsection{Balance equations}

This subsection presents the state-transition diagrams of Model-$X$ on both state spaces $S$ and $\widetilde{S}$, and derives the corresponding sets of balance equations.

\subsubsection{Balance equations for state space \texorpdfstring{$S$}{S}} \label{sec:gen:balance_equations}

The state-transition diagram of Model-$X$ on $S$ is given in Figure~\ref{fig:diagram-gen}.
\begin{figure}[H]
  \centering
  \resizebox{0.8\textwidth}{!}{%
    \begin{tikzpicture}[>=Stealth, line cap=round, line join=round]
  % --- global axes ---
  \draw[->] (0,0) -- (8,0) node[below right] {$n_1$};
  \draw[->] (0,0) -- (0,8) node[above left] {$n_2$};
  
  % Origin label
  \node[above right, xshift=2pt, yshift=2pt] at (0,0) {$(0,0)$};

  % --- Entry transitions from the origin ---
  \draw[->] (0,0) -- ++(0:1.2)  node[below] {$\lambda_1$};
  \draw[->] (0,0) -- ++(90:1.2) node[right] {$\lambda_2$};

  % --- boundary state (n1,0) ---
  \coordinate (A) at (6,0);
  \node[below] at (A) {$(n_1,0)$};
  \draw[->] (A) -- ++(0:1.2)   node[below] {$\lambda_1$}; 
  \draw[->] (A) -- ++(90:1.2)  node[right] {$\lambda_2$}; 
  \draw[->] (A) -- ++(180:1.2) node[below] {$\mu$};      
  \draw[->, dashed] (A) -- ++(135:1.2) node[left]  {$\gamma_1 n_1$};
  % Abandonment of class-1: parallel dotted arrow above mu
  \draw[->, dotted] (6, 0.3) -- (4.8, 0.3) node[above] {$\theta_1 n_1$};

  % --- interior state (n1,n2) ---
  \coordinate (B) at (6,4);
  \node[below left] at (B) {$(n_1,n_2)$};
  \draw[->] (B) -- ++(0:1.2)   node[below] {$\lambda_1$};
  \draw[->] (B) -- ++(90:1.2)  node[right] {$\lambda_2$};
  \draw[->] (B) -- ++(180:1.2) node[left]  {$\mu$};
  \draw[->, dashed] (B) -- ++(135:1.2) node[left]  {$\gamma_1 n_1$};
  \draw[->, dashed] (B) -- ++(315:1.2) node[right] {$\gamma_2 n_2$};
  % Abandonments: parallel dotted arrows
  \draw[->, dotted] (6, 4.3) -- (4.8, 4.3) node[above] {$\theta_1 n_1$};
  \draw[->, dotted] (B)      -- ++(270:1.2) node[left]  {$\theta_2 n_2$};

  % --- boundary state (0,n2) ---
  \coordinate (C) at (0,4);
  \node[left] at (C) {$(0,n_2)$};
  \draw[->] (C) -- ++(0:1.2)   node[below] {$\lambda_1$};
  \draw[->] (C) -- ++(90:1.2)  node[right] {$\lambda_2$};
  \draw[->] (C) -- ++(-90:1.2) node[left]  {$\mu$};  
  \draw[->, dashed] (C) -- ++(315:1.2) node[right] {$\gamma_2 n_2$};
  % Abandonment of class-2: parallel dotted arrow to the right of mu
  \draw[->, dotted] (0.3, 4) -- (0.3, 2.8) node[right] {$\theta_2 n_2$};

  % --- "idle" arcs ---
  \coordinate (idle) at (-1.25,-1.15);
  \node[below] at (idle) {$\text{idle}$};
  \draw[->] (idle) to[bend left=45] node[midway, above left] {$\lambda_1+\lambda_2$} (0,0);
  \draw[->] (0,0) to[bend left=45] node[midway, below right] {$\mu$} (idle);
\end{tikzpicture}
  }
  \caption{State transition diagram on state space $S$ for Model-$X$, with jockeying ($\gamma_i n_i$, dashed) and abandonments ($\theta_i n_i$, dotted).}
  \label{fig:diagram-gen}
\end{figure}

Next, we give the balance equations for the five parts of the state space: the interior ($n_1,n_2\ge1$), the two boundaries ($n_1=0$ and $n_2=0$), the busy state with empty queues $(0,0)$, and the idle state $(0)$. Recall that each waiting customer carries its own independent jockeying and patience clocks, so in a state with $n_i$ waiting class-$i$ customers the aggregate jockeying and abandonment rates are the linear terms $\gamma_i n_i$ and $\theta_i n_i$.
\begin{align}
    \begin{split}
        &\left[\lambda_1 + \lambda_2 + \mu + (\gamma_1 + \theta_1) n_1 + (\gamma_2 + \theta_2) n_2\right] \pi(n_1,n_2) \\
        & \quad = \quad \mu \pi(n_1+1,n_2) \\
        & \quad \quad + \lambda_1 \pi(n_1-1,n_2) \\
        & \quad \quad + \lambda_2 \pi(n_1,n_2-1) \\
        & \quad \quad + \gamma_1 (n_1+1) \pi(n_1+1,n_2-1) \\
        & \quad \quad + \gamma_2 (n_2+1) \pi(n_1-1,n_2+1) \\
        & \quad \quad + \theta_1 (n_1+1) \pi(n_1+1,n_2) \\
        & \quad \quad + \theta_2 (n_2+1) \pi(n_1,n_2+1)
    \end{split}
    && \text{for} \quad n_1 \geq 1, n_2 \geq 1 \label{eq:gen:balance_interior}\\
    \begin{split}
        & \left[ \lambda_1 + \lambda_2 + \mu +  (\gamma_2 + \theta_2) n_2 \right] \pi(0,n_2) \\
        & \quad = \quad \lambda_2 \pi(0,n_2-1) \\
        & \quad \quad + \mu \pi(1,n_2) \\
        & \quad \quad + \mu \pi(0,n_2+1) \\
        & \quad \quad + \gamma_1 \pi(1, n_2-1) \\
        & \quad \quad + \theta_1 \pi(1, n_2) \\
        & \quad \quad + \theta_2(n_2+1) \pi(0,n_2+1)
    \end{split}
    && \text{for} \quad n_1=0, n_2 \geq 1 \label{eq:gen:balance_y_boundary}\\
    \begin{split}
        & \left[ \lambda_1 + \lambda_2 + \mu + (\gamma_1 + \theta_1) n_1 \right] \pi(n_1,0) \\
        & \quad = \quad \lambda_1 \pi(n_1-1,0) \\
        & \quad \quad + \mu \pi(n_1+1, 0) \\
        & \quad \quad + \gamma_2 \pi(n_1-1, 1) \\
        & \quad \quad + \theta_1(n_1+1) \pi(n_1+1, 0) \\
        & \quad \quad + \theta_2 \pi(n_1,1)
    \end{split}
    && \text{for} \quad n_1 \geq 1, n_2 = 0 \label{eq:gen:balance_x_boundary}\\
    \begin{split}
        \left[\lambda_1+\lambda_2\right]\pi(0,0) = (\mu+\theta_1)\pi(1,0) + (\mu+\theta_2)\pi(0,1)
    \end{split} && \label{eq:gen:balance_zero}\\
    \begin{split}
        (\lambda_1+\lambda_2) \pi_0 = \mu \pi(0,0)
    \end{split} && \label{eq:gen:balance_idle}
\end{align}
Note that the idle-state equation~\eqref{eq:gen:balance_idle} is precisely the cut identity~\eqref{eq:gen:pi00_general} of Theorem~\ref{thm:gen:boundary}; it is repeated here so that the balance set is complete.

\subsubsection{Balance equations for state space \texorpdfstring{$\widetilde{S}$}{S~}}

Recall the state space $\widetilde{S}$ defined in Equation~\eqref{eq:gen:state_space_Stilde}, whose states $(n_2,n)$ record the number of waiting class-$2$ customers $n_2$ and the total number of waiting customers $n=n_1+n_2$. For readability we denote its limiting probabilities by $\tpi(n_2,n)$, for $(n_2,n)\in \widetilde{S}$. Figure~\ref{fig:diagram-gen-stilde} shows the corresponding state transition diagram: jockeying conserves $n$ and moves only $n_2$ ---a horizontal move in the diagram--- whereas arrivals increase $n$ by one and services and abandonments decrease it by one.

\begin{figure}[H]
  \centering
  \resizebox{0.8\textwidth}{!}{%
    \input{figures/diagram_jockeying_and_abandonments.tikz}
  }
  \caption{State transition diagram on state space $\widetilde{S}$.}
  \label{fig:diagram-gen-stilde}
\end{figure}

The balance equations, after combining the equations where $\tpi_0$ shows up, are:
\begin{align}
    \begin{split}
        & \left[ \lambda_1 + \lambda_2 + \mu + (\gamma_1+\theta_1)(n-n_2) + (\gamma_2+\theta_2) n_2 \right] \tpi(n_2,n) \\
        & \quad = \quad \lambda_1 \tpi(n_2,n-1) \\
        & \quad \quad + \lambda_2 \tpi(n_2-1,n-1) \\
        & \quad \quad + \mu \tpi(n_2, n+1) \\
        & \quad \quad + \gamma_1 (n-n_2+1) \tpi(n_2-1, n) \\
        & \quad \quad + \gamma_2 (n_2+1) \tpi(n_2+1,n) \\
        & \quad \quad + \theta_1 (n+1-n_2) \tpi(n_2, n+1) \\
        & \quad \quad + \theta_2 (n_2+1) \tpi(n_2+1, n+1)
    \end{split}
    && \text{for} \quad 1 \leq n_2 \leq n-1, \; n \geq 2 \label{eq:gen:tilde_balance_interior} \\
    \begin{split}
        & \left[ \lambda_1 + \lambda_2 + \mu + (\gamma_2+\theta_2) n \right] \tpi(n,n) \\
        & \quad = \quad \lambda_2 \tpi(n-1,n-1) \\
        & \quad \quad + \mu \tpi(n+1,n+1) \\
        & \quad \quad + \mu \tpi(n,n+1) \\
        & \quad \quad + \gamma_1 \tpi(n-1, n) \\
        & \quad \quad + \theta_1 \tpi(n, n+1) \\
        & \quad \quad + \theta_2 (n+1) \tpi(n+1, n+1)
    \end{split}
    && \text{for} \quad n_2=n, \; n \geq 1 \label{eq:gen:tilde_balance_diagonal} \\
    \begin{split}
        & \left[ \lambda_1 + \lambda_2 + \mu + (\gamma_1+\theta_1) n \right] \tpi(0,n) \\
        & \quad = \quad \lambda_1 \tpi(0,n-1) \\
        & \quad \quad + \mu \tpi(0, n+1) \\
        & \quad \quad + \gamma_2 \tpi(1,n) \\
        & \quad \quad + \theta_1 (n+1) \tpi(0, n+1) \\
        & \quad \quad + \theta_2 \tpi(1, n+1)
    \end{split}
    && \text{for} \quad n_2=0, \; n \geq 1 \label{eq:gen:tilde_balance_y_boundary} \\
    \begin{split}
        & \left[ \lambda_1 + \lambda_2 \right] \tpi(0,0) \\
        & \quad = \quad (\mu + \theta_1) \tpi(0,1) + (\mu + \theta_2) \tpi(1,1)
    \end{split}
    && \label{eq:gen:tilde_balance_zero} \\
    \begin{split}
        (\lambda_1+\lambda_2) \tpi_0 = \mu \tpi(0,0)
    \end{split}
    && \label{eq:gen:tilde_balance_idle}
\end{align}
As on $S$, the idle-state equation~\eqref{eq:gen:tilde_balance_idle} coincides with the cut identity~\eqref{eq:gen:pi00_general} of Theorem~\ref{thm:gen:boundary}.


\subsection{Equations for the (partial) probability generating functions} \label{ssec:pgf_derivation}

To study the limiting behaviour of this two-class system, we define the bivariate Probability Generating Function (PGF) of the joint queue lengths $N_1$ and $N_2$:
\begin{equation*}
    P(x,y) = \mathbb{E}[x^{N_1} y^{N_2}] = \sum_{n_1=0}^{\infty} \sum_{n_2=0}^{\infty} \pi(n_1,n_2) x^{n_1} y^{n_2}.
\end{equation*}

The derivation also uses the boundary terms, where at least one queue is empty:
\begin{align*}
    P_x(x) &= P(x,0) = \sum_{n_1=0}^{\infty} \pi(n_1,0) x^{n_1} \\
    P_y(y) &= P(0,y) = \sum_{n_2=0}^{\infty} \pi(0,n_2) y^{n_2}.
\end{align*}

Note that both terms include the busy-but-empty state $\pi(0,0)$, whereas the fully idle state is treated separately.

\subsubsection{Equation for PGF in state space \texorpdfstring{$S$}{S}}

\begin{lemma}[Fundamental equation for the PGF on $S$]\label{lem:gen:fundamental_equation}
Consider the two-class non-preemptive priority queue on the state space
$S=\{(n_1,n_2):n_1,n_2\ge 0\}$, with Poisson arrival rates $\lambda_1,\lambda_2$,
exponential service rate $\mu$, jockeying rates $\gamma_1,\gamma_2$ and abandonment
rates $\theta_1,\theta_2$, where $n_1$ and $n_2$ count the class-$1$ and class-$2$
customers \emph{waiting} in the queues. Assume the associated CTMC is positive
recurrent, so that the stationary distribution $\{\pi(n_1,n_2)\}$ exists and the
joint PGF $P(x,y)$ defined above is analytic on the closed unit polydisc
$|x|\le 1,\ |y|\le 1$. Then $P(x,y)$, together with the boundary function
$P_y(y)=P(0,y)$, satisfies the first-order linear partial differential equation
\begin{equation}
\begin{split}
    & \left[ (\lambda_1+\lambda_2+\mu)xy - \mu y - \lambda_1 x^2 y - \lambda_2 x y^2 \right] P(x,y) \\
    & + xy \left[ \gamma_1(x-y) + \theta_1(x-1) \right] \frac{\partial}{\partial x} P(x,y) \\
    & + xy \left[ \gamma_2(y-x) + \theta_2(y-1) \right] \frac{\partial}{\partial y} P(x,y) \\
    & \quad = \quad \mu(x-y) P_y(y) - \mu x(1-y)\,\pi(0,0).
\end{split}
\label{eq:gen:fundamental_equation}
\end{equation}
\end{lemma}

The structure of Equation~\eqref{eq:gen:fundamental_equation} is worth reading term by term, since it dictates the entire analysis that follows. The bracketed algebraic coefficient multiplying $P(x,y)$ is exactly the kernel of the baseline Model-$A$: the flux balance between the arrival contributions $\lambda_1 x^2 y$ and $\lambda_2 x y^2$, the service contribution $\mu y$, and the total outflow $(\lambda_1+\lambda_2+\mu)xy$.

The coefficients of the two partial derivatives encode the added mechanisms. Consider the coefficient $xy\,[\gamma_1(x-y)+\theta_1(x-1)]$ of the derivative in $x$. Its jockeying factor $(x-y)$ vanishes on the diagonal $x=y$: the analytic signature that jockeying merely relocates a customer and so conserves the total count. Its abandonment factor $(x-1)$ vanishes at $x=1$: the signature that abandonment respects normalisation ($P(1,1)=1$) yet removes a customer and thereby destroys conservation. The coefficient $xy\,[\gamma_2(y-x)+\theta_2(y-1)]$ of the derivative in $y$ states the symmetric facts for the class-$2$ mechanisms, with $(y-x)$ vanishing on the diagonal and $(y-1)$ at $y=1$.

The right-hand side involves the two boundary unknowns that any solution must determine: the boundary function $P_y(y)=P(0,y)$, the generating function of the states with an empty priority queue, and the scalar $\pi(0,0)$, the probability that the server is busy with both queues empty.

The position of these vanishing factors fixes, depending on the model, where the operative singularity lies and hence which technique applies. With both derivative terms absent (Model-$A$, \S\ref{sec:model_a}) the equation is algebraic and $P_y(y)$ is recovered from the kernel root $x^*(y)$ inside the unit disc. With one-way jockeying alone (Model-$B_2$, \S\ref{sec:model_b2}) the factor $(x-y)$ places the singularity at the interior point $x=y$ with a negative exponent, and $P_y(y)$ is determined by analyticity there. With class-$1$ abandonment alone (Model-$C_2$, \S\ref{sec:model_c2}) the factor $(x-1)$ places it on the boundary $x=1$ with a positive exponent, and $P_y(y)$ is determined by mere boundedness. Restricting a mechanism to the head of the queue removes the derivative term, reducing the equation to the algebraic Model-$A$ form (Models~$\BH$ and $\CH$, \S\ref{sec:model_b21}). The mirrored variants Model-$B_1$ (jockeying $2\to1$ only) and Model-$C_1$ (class-$2$ abandonment only), not treated in this thesis, place the derivative in $y$ instead; as for Model-$B$, the equation then couples the unknown boundary function $P_y(y)$ to the PGF along the integration path and reduces to a Volterra-type integral equation (\S\ref{sec:model_b}). These regimes are gathered side by side in Table~\ref{tab:comparison} of \S\ref{sec:comparison}, and are synthesised, together with the two irreducible obstructions, as the obstruction taxonomy of \S\ref{sec:conclusion}.

\begin{proof}[Proof of Lemma~\ref{lem:gen:fundamental_equation}]

    We multiply each of the balance equations~\eqref{eq:gen:balance_interior}--\eqref{eq:gen:balance_zero} by $x^{n_1}y^{n_2}$ and sum over its domain. Two intermediate transforms arise along the way, describing the boundary dynamics when one queue holds exactly one waiting customer:
    \begin{equation*}
        P_y(y \,|\, n_1{=}1) := \sum_{n_2=0}^{\infty} \pi(1,n_2)\, y^{n_2},
        \qquad
        P_x(x \,|\, n_2{=}1) := \sum_{n_1=0}^{\infty} \pi(n_1,1)\, x^{n_1}.
    \end{equation*}
    Despite the conditional-looking bar, these are joint (unnormalised) transforms of the states with exactly one waiting class-$1$ (respectively class-$2$) customer; the derivation eliminates them in favour of the primary unknowns $P_x(x)$, $P_y(y)$ and $\pi(0,0)$.

    From the interior equation~\eqref{eq:gen:balance_interior}:
    \begin{align*}
        \begin{split}
            & (\lambda_1 + \lambda_2 + \mu) \sum_{n_1=1}^{\infty} \sum_{n_2=1}^{\infty} \pi(n_1,n_2) x^{n_1}y^{n_2} \\
            &+ \gamma_1 \sum_{n_1=1}^{\infty} \sum_{n_2=1}^{\infty} n_1 \pi(n_1,n_2) x^{n_1}y^{n_2} \\
            &+ \gamma_2 \sum_{n_1=1}^{\infty} \sum_{n_2=1}^{\infty} n_2 \pi(n_1,n_2) x^{n_1}y^{n_2} \\
            &+ \theta_1 \sum_{n_1=1}^{\infty} \sum_{n_2=1}^{\infty} n_1 \pi(n_1,n_2) x^{n_1}y^{n_2} \\
            &+ \theta_2 \sum_{n_1=1}^{\infty} \sum_{n_2=1}^{\infty} n_2 \pi(n_1,n_2) x^{n_1}y^{n_2} \\
            &\quad = \quad \mu \sum_{n_1=1}^{\infty} \sum_{n_2=1}^{\infty} \pi(n_1+1,n_2)x^{n_1}y^{n_2} \\
            &\quad\quad  + \lambda_1 \sum_{n_1=1}^{\infty} \sum_{n_2=1}^{\infty} \pi(n_1-1,n_2)x^{n_1}y^{n_2} \\
            &\quad\quad + \lambda_2 \sum_{n_1=1}^{\infty} \sum_{n_2=1}^{\infty} \pi(n_1,n_2-1)x^{n_1}y^{n_2} \\
            &\quad\quad + \gamma_1 \sum_{n_1=1}^{\infty} \sum_{n_2=1}^{\infty} (n_1+1)\pi(n_1+1,n_2-1)x^{n_1}y^{n_2} \\
            &\quad\quad + \gamma_2 \sum_{n_1=1}^{\infty} \sum_{n_2=1}^{\infty} (n_2+1)\pi(n_1-1,n_2+1)x^{n_1}y^{n_2} \\
            &\quad\quad + \theta_1 \sum_{n_1=1}^{\infty} \sum_{n_2=1}^{\infty} (n_1+1)\pi(n_1+1,n_2)x^{n_1}y^{n_2} \\
            &\quad\quad + \theta_2 \sum_{n_1=1}^{\infty} \sum_{n_2=1}^{\infty} (n_2+1)\pi(n_1,n_2+1)x^{n_1}y^{n_2}
        \end{split} \\
        \begin{split}
            & LHS_1 + LHS_2 + LHS_3 + LHS_4 + LHS_5 \\
            & \quad = \quad RHS_1 + RHS_2 + RHS_3 + RHS_4 + RHS_5 + RHS_6 + RHS_7
        \end{split}
    \end{align*}
    where $LHS_i$ and $RHS_j$ denote the $i$-th and $j$-th summands on the left- and right-hand sides of the display above, in the order written. Evaluating each summand:
    \begin{align*}
        \begin{split}
            & LHS_1 \\
            & \quad = \quad (\lambda_1 + \lambda_2 + \mu) \sum_{n_1=1}^{\infty} \sum_{n_2=1}^{\infty} \pi(n_1,n_2)x^{n_1}y^{n_2} \\
            & \quad = \quad (\lambda_1 + \lambda_2 + \mu) \sum_{n_1=1}^{\infty} \left[ \sum_{n_2=0}^{\infty} \pi(n_1,n_2) x^{n_1}y^{n_2} - \pi(n_1,0) x^{n_1} y^0 \right] \\
            & \quad = \quad (\lambda_1 + \lambda_2 + \mu) \left[
                \sum_{n_1=1}^{\infty} \sum_{n_2=0}^{\infty} \pi(n_1,n_2) x^{n_1}y^{n_2}
                - \sum_{n_1=1}^{\infty} \pi(n_1,0) x^{n_1}
            \right] \\
            & \quad = \quad (\lambda_1 + \lambda_2 + \mu) \left[
                \sum_{n_1=0}^{\infty} \sum_{n_2=0}^{\infty} \pi(n_1,n_2) x^{n_1}y^{n_2}
                - \sum_{n_2=0}^{\infty} \pi(0,n_2) y^{n_2}
            \right] \\
            & \quad \quad - \; (\lambda_1 + \lambda_2 + \mu) \left[
                \sum_{n_1=0}^{\infty} \pi(n_1,0) x^{n_1} - \pi(0,0)
            \right] \\
            & \quad = \quad (\lambda_1 + \lambda_2 + \mu) \left[ P(x,y) - P_y(y) - P_x(x) + \pi(0,0) \right]
        \end{split} \\
        \begin{split}
            & LHS_2 \\
            & \quad = \quad \gamma_1 \sum_{n_1=1}^{\infty} \sum_{n_2=1}^{\infty} n_1 \pi(n_1,n_2) x^{n_1} y^{n_2} \\
            & \quad = \quad \gamma_1 x \sum_{n_1=1}^{\infty} \sum_{n_2=1}^{\infty} n_1 \pi(n_1,n_2) x^{n_1-1} y^{n_2} \\
            & \quad = \quad \gamma_1 x \left[
                \sum_{n_1=1}^{\infty} \sum_{n_2=0}^{\infty} n_1 \pi(n_1,n_2) x^{n_1-1} y^{n_2}
                - \sum_{n_1=1}^{\infty} n_1 \pi(n_1,0) x^{n_1-1}
            \right] \\
            & \quad = \quad \gamma_1 x \left[
                \sum_{n_1=0}^{\infty} \sum_{n_2=0}^{\infty} n_1 \pi(n_1,n_2) x^{n_1-1} y^{n_2}
                - \sum_{n_1=0}^{\infty} n_1 \pi(n_1,0) x^{n_1-1}
            \right] \\
            & \quad = \quad \gamma_1 x \left[
                \frac{\partial}{\partial x} P(x,y) - \frac{d}{dx} P_x(x)
            \right]
        \end{split} \\
        \begin{split}
            & LHS_3 \\
            & \quad = \quad \gamma_2 \sum_{n_1=1}^{\infty} \sum_{n_2=1}^{\infty} n_2 \pi(n_1,n_2) x^{n_1} y^{n_2} \\
            & \quad = \quad \gamma_2 \sum_{n_1=1}^{\infty} \left[ \sum_{n_2=0}^{\infty} n_2 \pi(n_1,n_2) x^{n_1} y^{n_2} - 0 \cdot \pi(n_1,0) x^{n_1} \right] \\
            & \quad = \quad \gamma_2 y \left[
                \sum_{n_1=0}^{\infty} \sum_{n_2=0}^{\infty} n_2 \pi(n_1,n_2) x^{n_1} y^{n_2-1}
                - \sum_{n_2=0}^{\infty} n_2 \pi(0,n_2) y^{n_2-1}
            \right] \\
            & \quad = \quad \gamma_2 y \left[
                \frac{\partial}{\partial y} P(x,y) - \frac{d}{dy} P_y(y)
            \right]
        \end{split} \\
        \begin{split}
            & LHS_4 \\
            & \quad = \quad \theta_1 \sum_{n_1=1}^{\infty} \sum_{n_2=1}^{\infty} n_1 \pi(n_1,n_2) x^{n_1} y^{n_2} \\
            & \quad = \quad \theta_1 x \left[
                \frac{\partial}{\partial x} P(x,y) - \frac{d}{dx} P_x(x)
            \right]
            \qquad \text{(the computation of $LHS_2$ with $\theta_1$ in place of $\gamma_1$)}
        \end{split} \\
        \begin{split}
            & LHS_5 \\
            & \quad = \quad \theta_2 \sum_{n_1=1}^{\infty} \sum_{n_2=1}^{\infty} n_2 \pi(n_1,n_2) x^{n_1} y^{n_2} \\
            & \quad = \quad \theta_2 y \left[
                \frac{\partial}{\partial y} P(x,y) - \frac{d}{dy} P_y(y)
            \right]
            \qquad \text{(the computation of $LHS_3$ with $\theta_2$ in place of $\gamma_2$)}
        \end{split} \\
        \begin{split}
            & RHS_1 \\
            & \quad = \quad \mu \sum_{n_1=1}^{\infty} \sum_{n_2=1}^{\infty} \pi(n_1+1, n_2) x^{n_1} y^{n_2} \\
            & \quad = \quad \mu x^{-1} \sum_{n_1=1}^{\infty} \sum_{n_2=1}^{\infty} \pi(n_1+1, n_2) x^{n_1+1} y^{n_2} \\
            & \quad = \quad \mu x^{-1} \sum_{m=2}^{\infty} \sum_{n_2=1}^{\infty} \pi(m, n_2) x^{m} y^{n_2} \\
            & \quad = \quad \mu x^{-1} \sum_{m=2}^{\infty} \left[
                \sum_{n_2=0}^{\infty} \pi(m, n_2) x^{m} y^{n_2} - \pi(m, 0) x^m
            \right] \\
            & \quad = \quad \mu x^{-1} \left[ \sum_{m=2}^{\infty} \sum_{n_2=0}^{\infty} \pi(m, n_2) x^{m} y^{n_2}
                - \sum_{m=2}^{\infty} \pi(m, 0) x^{m}
            \right] \\
            & \quad = \quad \mu x^{-1} \left[
                \sum_{m=0}^{\infty} \sum_{n_2=0}^{\infty} \pi(m, n_2) x^{m} y^{n_2}
                - x \sum_{n_2=0}^{\infty} \pi(1, n_2) y^{n_2}
                - \sum_{n_2=0}^{\infty} \pi(0, n_2) y^{n_2}
            \right] \\
            & \quad \quad - \; \mu x^{-1} \left[
                \sum_{m=0}^{\infty} \pi(m, 0) x^{m} - x \pi(1, 0) - \pi(0, 0)
            \right] \\
            & \quad = \quad \mu x^{-1} \left[
                P(x,y) - x P_y(y | n_1=1) - P_y(y) - P_x(x) + x \pi(1, 0) + \pi(0, 0)
            \right]
        \end{split} \\
        \begin{split}
            & RHS_2 \\
            & \quad = \quad \lambda_1 \sum_{n_1=1}^{\infty} \sum_{n_2=1}^{\infty} \pi(n_1-1, n_2) x^{n_1} y^{n_2} \\
            & \quad = \quad \lambda_1 x \sum_{n_1=1}^{\infty} \sum_{n_2=1}^{\infty} \pi(n_1-1, n_2) x^{n_1-1} y^{n_2} \\
            & \quad = \quad \lambda_1 x \sum_{m=0}^{\infty} \sum_{n_2=1}^{\infty} \pi(m, n_2) x^{m} y^{n_2} \\
            & \quad = \quad \lambda_1 x \sum_{m=0}^{\infty} \left[
                \sum_{n_2=0}^{\infty} \pi(m, n_2) x^{m} y^{n_2} - \pi(m, 0) x^m
            \right] \\
            & \quad = \quad \lambda_1 x \left[
                \sum_{m=0}^{\infty} \sum_{n_2=0}^{\infty} \pi(m, n_2) x^{m} y^{n_2}
                - \sum_{m=0}^{\infty} \pi(m, 0) x^m
            \right] \\
            & \quad = \quad \lambda_1 x \left[ P(x,y) - P_x(x) \right]
        \end{split} \\
        \begin{split}
            & RHS_3 \\
            & \quad = \quad \lambda_2 \sum_{n_1=1}^{\infty} \sum_{n_2=1}^{\infty} \pi(n_1, n_2-1) x^{n_1} y^{n_2} \\
            & \quad = \quad \lambda_2 y \sum_{n_1=1}^{\infty} \sum_{n_2=1}^{\infty} \pi(n_1, n_2-1) x^{n_1} y^{n_2-1} \\
            & \quad = \quad \lambda_2 y \sum_{n_1=1}^{\infty} \sum_{m=0}^{\infty} \pi(n_1, m) x^{n_1} y^{m} \\
            & \quad = \quad \lambda_2 y \left[
                \sum_{n_1=0}^{\infty} \sum_{m=0}^{\infty} \pi(n_1, m) x^{n_1} y^{m} - \sum_{m=0}^{\infty} \pi(0, m) y^m
            \right] \\
            & \quad = \quad \lambda_2 y \left[ P(x,y) - P_y(y) \right]
        \end{split} \\
        \begin{split}
            & RHS_4 \\
            & \quad = \quad \gamma_1 \sum_{n_1=1}^{\infty} \sum_{n_2=1}^{\infty} (n_1+1) \pi(n_1+1, n_2-1) x^{n_1} y^{n_2} \\
            & \quad = \quad \gamma_1 y \sum_{n_1=1}^{\infty} \sum_{m_2=0}^{\infty} (n_1+1) \pi(n_1+1, m_2) x^{n_1} y^{m_2} \\
            & \quad = \quad \gamma_1 y \sum_{m_1=2}^{\infty} \sum_{m_2=0}^{\infty} m_1 \pi(m_1, m_2) x^{m_1-1} y^{m_2} \\
            & \quad = \quad \gamma_1 y \left[
                \sum_{m_1=0}^{\infty} \sum_{m_2=0}^{\infty} m_1 \pi(m_1, m_2) x^{m_1-1} y^{m_2}
                - \sum_{m_2=0}^{\infty} \pi(1, m_2) y^{m_2}
            \right] \\
            & \quad = \quad \gamma_1 y \left[
                \frac{\partial}{\partial x} P(x,y) - P_y(y | n_1=1)
            \right]
        \end{split} \\
        \begin{split}
            & RHS_5 \\
            & \quad = \quad \gamma_2 \sum_{n_1=1}^{\infty} \sum_{n_2=1}^{\infty} (n_2+1) \pi(n_1-1, n_2+1) x^{n_1} y^{n_2} \\
            & \quad = \quad \gamma_2 \sum_{n_1=1}^{\infty} \sum_{m_2=2}^{\infty} m_2 \pi(n_1-1, m_2) x^{n_1} y^{m_2-1} \\
            & \quad = \quad \gamma_2 x \left[
                \sum_{n_1=1}^{\infty} \sum_{m_2=0}^{\infty} m_2 \pi(n_1-1, m_2) x^{n_1-1} y^{m_2-1}
                - \sum_{n_1=1}^{\infty} \pi(n_1-1, 1) x^{n_1-1}
            \right] \\
            & \quad = \quad \gamma_2 x \left[
                \sum_{m_1=0}^{\infty} \sum_{m_2=0}^{\infty} m_2 \pi(m_1, m_2) x^{m_1} y^{m_2-1}
                - \sum_{m_1=0}^{\infty} \pi(m_1, 1) x^{m_1}
            \right] \\
            & \quad = \quad \gamma_2 x \left[
                \frac{\partial}{\partial y} P(x,y) - P_x(x | n_2=1)
            \right]
        \end{split} \\
        \begin{split}
            & RHS_6 \\
            & \quad = \quad \theta_1 \sum_{n_1=1}^{\infty} \sum_{n_2=1}^{\infty} (n_1+1) \pi(n_1+1, n_2) x^{n_1} y^{n_2} \\
            & \quad = \quad \theta_1 \sum_{m_1=2}^{\infty} \sum_{n_2=1}^{\infty} m_1 \pi(m_1, n_2) x^{m_1-1} y^{n_2} \\
            & \quad = \quad \theta_1 \left[
                \sum_{m_1=0}^{\infty} \sum_{n_2=1}^{\infty} m_1 \pi(m_1, n_2) x^{m_1-1} y^{n_2}
                - \sum_{n_2=1}^{\infty} \pi(1, n_2) y^{n_2}
            \right] \\
            & \quad = \quad \theta_1 \left[
                \sum_{m_1=0}^{\infty} \sum_{n_2=0}^{\infty} m_1 \pi(m_1, n_2) x^{m_1-1} y^{n_2}
                - \sum_{m_1=0}^{\infty} m_1 \pi(m_1, 0) x^{m_1-1}
            \right] \\
            & \quad \quad - \theta_1 \left[
                \sum_{n_2=0}^{\infty} \pi(1, n_2) y^{n_2}
                - \pi(1, 0)
            \right] \\
            & \quad = \quad \theta_1 \left[
                \frac{\partial}{\partial x} P(x,y) - \frac{d}{dx} P_x(x) - P_y(y | n_1=1) + \pi(1, 0)
            \right]
        \end{split} \\
        \begin{split}
            & RHS_7 \\
            & \quad = \quad \theta_2 \sum_{n_1=1}^{\infty} \sum_{n_2=1}^{\infty} (n_2+1) \pi(n_1, n_2+1) x^{n_1} y^{n_2} \\
            & \quad = \quad \theta_2 \sum_{n_1=1}^{\infty} \sum_{m_2=2}^{\infty} m_2 \pi(n_1, m_2) x^{n_1} y^{m_2-1} \\
            & \quad = \quad \theta_2 \left[
                \sum_{n_1=1}^{\infty} \sum_{m_2=0}^{\infty} m_2 \pi(n_1, m_2) x^{n_1} y^{m_2-1}
                - \sum_{n_1=1}^{\infty} \pi(n_1, 1) x^{n_1}
            \right] \\
            & \quad = \quad \theta_2 \left[
                \sum_{n_1=0}^{\infty} \sum_{m_2=0}^{\infty} m_2 \pi(n_1, m_2) x^{n_1} y^{m_2-1}
                - \sum_{m_2=0}^{\infty} m_2 \pi(0, m_2) y^{m_2-1}
            \right] \\
            & \quad \quad - \theta_2 \left[
                \sum_{n_1=0}^{\infty} \pi(n_1, 1) x^{n_1}
                - \pi(0, 1)
            \right] \\
            & \quad = \quad \theta_2 \left[
                \frac{\partial}{\partial y} P(x,y) - \frac{d}{dy} P_y(y) - P_x(x | n_2=1) + \pi(0, 1)
            \right]
        \end{split}
    \end{align*}
    
    Substituting the evaluated summands back into the transformed balance equation:
    \begin{align}
        \begin{split}
            & (\lambda_1 + \lambda_2 + \mu) \left[
                P(x,y) - P_y(y) - P_x(x) + \pi(0,0)
            \right] \\
            & + \gamma_1 x \left[ \frac{\partial}{\partial x} P(x,y) - \frac{d}{dx} P_x(x) \right] \\
            & + \gamma_2 y \left[ \frac{\partial}{\partial y} P(x,y) - \frac{d}{dy} P_y(y) \right] \\
            & + \theta_1 x \left[ \frac{\partial}{\partial x} P(x,y) - \frac{d}{dx} P_x(x) \right] \\
            & + \theta_2 y \left[ \frac{\partial}{\partial y} P(x,y) - \frac{d}{dy} P_y(y) \right] \\
            & \quad = \quad \mu x^{-1} \left[
                P(x,y) - x P_y(y | n_1=1) - P_y(y) - P_x(x) + x \pi(1,0) + \pi(0,0)
            \right] \\
            & \quad \quad + \lambda_1 x \left[ P(x,y) - P_x(x) \right] \\
            & \quad \quad + \lambda_2 y \left[ P(x,y) - P_y(y) \right] \\
            & \quad \quad + \gamma_1 y \left[ \frac{\partial}{\partial x} P(x,y) - P_y(y | n_1=1) \right] \\
            & \quad \quad + \gamma_2 x \left[ \frac{\partial}{\partial y} P(x,y) - P_x(x | n_2=1) \right] \\
            & \quad \quad + \theta_1 \left[
                \frac{\partial}{\partial x} P(x,y) - \frac{d}{dx} P_x(x) - P_y(y | n_1=1) + \pi(1, 0)
            \right] \\
            & \quad \quad + \theta_2 \left[
                \frac{\partial}{\partial y} P(x,y) - \frac{d}{dy} P_y(y) - P_x(x | n_2=1) + \pi(0, 1)
            \right]
        \end{split} \nonumber \\
        \intertext{Grouping the coefficients of each transform yields}
        \begin{split}
            & \left[ (\lambda_1+\lambda_2+\mu) - \mu x^{-1} - \lambda_1 x - \lambda_2 y \right] P(x,y) \\
            & + \left[ \gamma_1(x-y) + \theta_1(x-1) \right] \frac{\partial}{\partial x} P(x,y) \\
            & + \left[ \gamma_2(y-x) + \theta_2(y-1) \right] \frac{\partial}{\partial y} P(x,y) \\
            & \quad = \quad \left[ (\lambda_1 + \lambda_2 + \mu) - \mu x^{-1} - \lambda_1 x \right] P_x(x) \\
            & \quad \quad + \left[ (\lambda_1 + \lambda_2 + \mu) - \mu x^{-1} - \lambda_2 y \right] P_y(y) \\
            & \quad \quad + \left[ -(\lambda_1 + \lambda_2 + \mu) + \mu x^{-1} \right] \pi(0,0) \\
            & \quad \quad + \left[ \gamma_1 x + \theta_1(x-1) \right] \frac{d}{dx} P_x(x) \\
            & \quad \quad + \left[ \gamma_2 y + \theta_2(y-1) \right] \frac{d}{dy} P_y(y) \\
            & \quad \quad - \left[ \mu + \gamma_1 y + \theta_1 \right] P_y(y | n_1=1) \\
            & \quad \quad - \left[ \gamma_2 x + \theta_2 \right] P_x(x | n_2=1) \\
            & \quad \quad + \mu \pi(1,0) \\
            & \quad \quad + \theta_1 \pi(1,0) \\
            & \quad \quad + \theta_2 \pi(0,1)
        \end{split} \nonumber \\
        \intertext{and multiplying through by $xy$ to clear the $x^{-1}$ factors:}
        \begin{split}
            & \left[ (\lambda_1+\lambda_2+\mu)xy - \mu y - \lambda_1 x^2 y - \lambda_2 x y^2 \right] P(x,y) \\
            & + xy \left[ \gamma_1(x-y) + \theta_1(x-1) \right] \frac{\partial}{\partial x} P(x,y) \\
            & + xy \left[ \gamma_2(y-x) + \theta_2(y-1) \right] \frac{\partial}{\partial y} P(x,y) \\
            & \quad = \quad \left[ (\lambda_1 + \lambda_2 + \mu)xy - \mu y - \lambda_1 x^2 y \right] P_x(x) \\
            & \quad \quad + \left[ (\lambda_1 + \lambda_2 + \mu)xy - \mu y - \lambda_2 x y^2 \right] P_y(y) \\
            & \quad \quad - \left[ (\lambda_1 + \lambda_2 + \mu)xy - \mu y \right] \pi(0,0) \\
            & \quad \quad + xy \left[ \gamma_1 x + \theta_1(x-1) \right] \frac{d}{dx} P_x(x) \\
            & \quad \quad + xy \left[ \gamma_2 y + \theta_2(y-1) \right] \frac{d}{dy} P_y(y) \\
            & \quad \quad - xy \left[ \mu + \gamma_1 y + \theta_1 \right] P_y(y | n_1=1) \\
            & \quad \quad - xy \left[ \gamma_2 x + \theta_2 \right] P_x(x | n_2=1) \\
            & \quad \quad + (\mu + \theta_1) xy \, \pi(1,0) \\
            & \quad \quad + \theta_2 xy \, \pi(0,1)
        \end{split} \label{eq:gen:transform_interior}
    \end{align}
    
    The $x$-boundary equation~\eqref{eq:gen:balance_x_boundary}, treated analogously, gives:
    \begin{align}
        \begin{split}
            & \sum_{n_1=1}^{\infty} \left[ \lambda_1 + \lambda_2 + \mu + (\gamma_1 + \theta_1) n_1 \right] \pi(n_1, 0) x^{n_1} \\
            & \quad = \quad \lambda_1 \sum_{n_1=1}^{\infty} \pi(n_1-1,0) x^{n_1}
            + \mu \sum_{n_1=1}^{\infty} \pi(n_1+1,0) x^{n_1} \\
            & \quad \quad + \gamma_2 \sum_{n_1=1}^{\infty} \pi(n_1-1, 1) x^{n_1}
            + \theta_1 \sum_{n_1=1}^{\infty} (n_1+1) \pi(n_1+1, 0) x^{n_1}
            + \theta_2 \sum_{n_1=1}^{\infty} \pi(n_1, 1) x^{n_1}
        \end{split} \nonumber \\
        \intertext{Extracting the rate prefactors and shifting the summation indices:}
        \begin{split}
            & (\lambda_1 + \lambda_2 + \mu) \sum_{n_1=1}^{\infty} \pi(n_1, 0) x^{n_1}
            + (\gamma_1 + \theta_1) x \sum_{n_1=1}^{\infty} n_1 \pi(n_1, 0) x^{n_1-1} \\
            & \quad = \quad \lambda_1 x \sum_{m=0}^{\infty} \pi(m, 0) x^{m}
            + \mu x^{-1} \sum_{m=2}^{\infty} \pi(m, 0) x^{m} \\
            & \quad \quad + \gamma_2 x \sum_{m=0}^{\infty} \pi(m, 1) x^{m}
            + \theta_1 \sum_{m=2}^{\infty} m \pi(m, 0) x^{m-1}
            + \theta_2 \sum_{m=1}^{\infty} \pi(m, 1) x^{m}
        \end{split} \nonumber \\
        \intertext{Completing each sum to start at zero and subtracting the added terms:}
        \begin{split}
            & (\lambda_1 + \lambda_2 + \mu) \left[ \sum_{n_1=0}^{\infty} \pi(n_1, 0) x^{n_1} - \pi(0,0) \right]
            + (\gamma_1 + \theta_1) x \sum_{n_1=0}^{\infty} n_1 \pi(n_1, 0) x^{n_1-1} \\
            & \quad = \quad \lambda_1 x \sum_{m=0}^{\infty} \pi(m, 0) x^{m} \\
            & \quad \quad + \mu x^{-1} \left[ \sum_{m=0}^{\infty} \pi(m, 0) x^{m} - x \pi(1, 0) - \pi(0,0) \right] \\
            & \quad \quad + \gamma_2 x \sum_{m=0}^{\infty} \pi(m, 1) x^{m} \\
            & \quad \quad + \theta_1 \left[ \sum_{m=0}^{\infty} m \pi(m, 0) x^{m-1} - \pi(1, 0) \right] \\
            & \quad \quad + \theta_2 \left[ \sum_{m=0}^{\infty} \pi(m, 1) x^{m} - \pi(0, 1) \right]
        \end{split} \nonumber \\
        \intertext{Recognising the transforms defined in \S\ref{ssec:pgf_derivation}:}
        \begin{split}
            & (\lambda_1 + \lambda_2 + \mu) \left[ P_x(x) - \pi(0,0) \right]
            + (\gamma_1 + \theta_1) x \frac{d}{dx} P_x(x) \\
            & \quad = \quad \lambda_1 x P_x(x)
            + \mu x^{-1} \left[ P_x(x) - x \pi(1, 0) - \pi(0, 0) \right] \\
            & \quad \quad + \gamma_2 x P_x(x | n_2=1)
            + \theta_1 \left[ \frac{d}{dx} P_x(x) - \pi(1, 0) \right]
            + \theta_2 \left[ P_x(x | n_2=1) - \pi(0, 1) \right]
        \end{split} \nonumber \\
        \intertext{Grouping the coefficients of $P_x(x)$ and of its derivative:}
        \begin{split}
            & \left[ (\lambda_1 + \lambda_2 + \mu) - \mu x^{-1} - \lambda_1 x \right] P_x(x)
            + \left[ \gamma_1 x + \theta_1(x-1) \right] \frac{d}{dx} P_x(x) \\
            & \quad = \quad \left[ (\lambda_1 + \lambda_2 + \mu) - \mu x^{-1} \right] \pi(0,0)
            - (\mu + \theta_1) \pi(1, 0)
            - \theta_2 \pi(0, 1) \\
            & \quad \quad + \left[ \gamma_2 x + \theta_2 \right] P_x(x | n_2=1)
        \end{split} \nonumber \\
        \intertext{and finally multiplying through by $xy$, ready for substitution into Equation~\eqref{eq:gen:transform_interior}:}
        \begin{split}
            & \left[ (\lambda_1 + \lambda_2 + \mu)xy - \mu y - \lambda_1 x^2 y \right] P_x(x)
            + xy \left[ \gamma_1 x + \theta_1(x-1) \right] \frac{d}{dx} P_x(x) \\
            & \quad = \quad \left[ (\lambda_1 + \lambda_2 + \mu)xy - \mu y \right] \pi(0,0) \\
            & \quad \quad - (\mu + \theta_1) xy \, \pi(1, 0) \\
            & \quad \quad - \theta_2 xy \, \pi(0, 1) \\
            & \quad \quad + xy \left[ \gamma_2 x + \theta_2 \right] P_x(x | n_2=1)
        \end{split} \label{eq:gen:transform_x_boundary}
    \end{align}
    
    The $y$-boundary equation~\eqref{eq:gen:balance_y_boundary} gives:
    \begin{align}
        \begin{split}
            & \sum_{n_2=1}^{\infty} \left[ \lambda_1 + \lambda_2 + \mu + (\gamma_2 + \theta_2) n_2 \right] \pi(0, n_2) y^{n_2} \\
            & \quad = \quad \lambda_2 \sum_{n_2=1}^{\infty} \pi(0, n_2-1) y^{n_2} \\
            & \quad \quad + \mu \sum_{n_2=1}^{\infty} \pi(1, n_2) y^{n_2} \\
            & \quad \quad + \mu \sum_{n_2=1}^{\infty} \pi(0, n_2+1) y^{n_2} \\
            & \quad \quad + \gamma_1 \sum_{n_2=1}^{\infty} \pi(1, n_2-1) y^{n_2} \\
            & \quad \quad + \theta_1 \sum_{n_2=1}^{\infty} \pi(1, n_2) y^{n_2} \\
            & \quad \quad + \theta_2 \sum_{n_2=1}^{\infty} (n_2+1) \pi(0, n_2+1) y^{n_2}
        \end{split} \nonumber \\
        \intertext{Extracting the rate prefactors and shifting the summation indices:}
        \begin{split}
            & (\lambda_1 + \lambda_2 + \mu) \sum_{n_2=1}^{\infty} \pi(0, n_2) y^{n_2} \\
            & + (\gamma_2 + \theta_2) y \sum_{n_2=1}^{\infty} n_2 \pi(0, n_2) y^{n_2-1} \\
            & \quad = \quad \lambda_2 y \sum_{m=0}^{\infty} \pi(0, m) y^{m} \\
            & \quad \quad + \mu \sum_{n_2=1}^{\infty} \pi(1, n_2) y^{n_2} \\
            & \quad \quad + \mu y^{-1} \sum_{m=2}^{\infty} \pi(0, m) y^{m} \\
            & \quad \quad + \gamma_1 y \sum_{m=0}^{\infty} \pi(1, m) y^{m} \\
            & \quad \quad + \theta_1 \sum_{n_2=1}^{\infty} \pi(1, n_2) y^{n_2} \\
            & \quad \quad + \theta_2 \sum_{m=2}^{\infty} m \pi(0, m) y^{m-1}
        \end{split} \nonumber \\
        \intertext{Completing each sum and subtracting the added terms:}
        \begin{split}
            & (\lambda_1 + \lambda_2 + \mu) \left[ \sum_{n_2=0}^{\infty} \pi(0, n_2) y^{n_2} - \pi(0,0) \right] \\
            & + (\gamma_2 + \theta_2) y \sum_{n_2=0}^{\infty} n_2 \pi(0, n_2) y^{n_2-1} \\
            & \quad = \quad \lambda_2 y \sum_{m=0}^{\infty} \pi(0, m) y^{m} \\
            & \quad \quad + \mu \left[ \sum_{n_2=0}^{\infty} \pi(1, n_2) y^{n_2} - \pi(1, 0) \right] \\
            & \quad \quad + \mu y^{-1} \left[ \sum_{m=0}^{\infty} \pi(0, m) y^{m} - y \pi(0, 1) - \pi(0, 0) \right] \\
            & \quad \quad + \gamma_1 y \sum_{m=0}^{\infty} \pi(1, m) y^{m} \\
            & \quad \quad + \theta_1 \left[ \sum_{n_2=0}^{\infty} \pi(1, n_2) y^{n_2} - \pi(1, 0) \right] \\
            & \quad \quad + \theta_2 \left[ \sum_{m=0}^{\infty} m \pi(0, m) y^{m-1} - \pi(0, 1) \right]
        \end{split} \nonumber \\
        \intertext{Recognising the transforms:}
        \begin{split}
            & (\lambda_1 + \lambda_2 + \mu) \left[ P_y(y) - \pi(0, 0) \right]
            + (\gamma_2 + \theta_2) y \frac{d}{dy} P_y(y) \\
            & \quad = \quad \lambda_2 y P_y(y)
            + \mu \left[ P_y(y | n_1=1) - \pi(1, 0) \right] \\
            & \quad \quad + \mu y^{-1} \left[ P_y(y) - y \pi(0, 1) - \pi(0, 0) \right]
            + \gamma_1 y P_y(y | n_1=1) \\
            & \quad \quad + \theta_1 \left[ P_y(y | n_1=1) - \pi(1, 0) \right]
            + \theta_2 \left[ \frac{d}{dy} P_y(y) - \pi(0, 1) \right]
        \end{split} \nonumber \\
        \intertext{Grouping the coefficients of $P_y(y)$ and of its derivative:}
        \begin{split}
            & \left[ (\lambda_1 + \lambda_2 + \mu) - \mu y^{-1} - \lambda_2 y \right] P_y(y)
            + \left[ \gamma_2 y + \theta_2(y-1) \right] \frac{d}{dy} P_y(y) \\
            & \quad = \quad \left[ \mu + \gamma_1 y + \theta_1 \right] P_y(y | n_1=1) \\
            & \quad \quad - (\mu + \theta_1) \pi(1, 0)
            - (\mu + \theta_2) \pi(0, 1) \\
            & \quad \quad + \left[ (\lambda_1 + \lambda_2 + \mu) - \mu y^{-1} \right] \pi(0, 0)
        \end{split} \nonumber \\
        \intertext{Substituting the empty-queues balance equation~\eqref{eq:gen:balance_zero} for $(\mu+\theta_1)\pi(1,0)+(\mu+\theta_2)\pi(0,1)$:}
        \begin{split}
            & \left[ (\lambda_1 + \lambda_2 + \mu) - \mu y^{-1} - \lambda_2 y \right] P_y(y)
            + \left[ \gamma_2 y + \theta_2(y-1) \right] \frac{d}{dy} P_y(y) \\
            & \quad = \quad \left[ \mu + \gamma_1 y + \theta_1 \right] P_y(y | n_1=1) \\
            & \quad \quad + \left[ (\lambda_1 + \lambda_2 + \mu) - \mu y^{-1} - (\lambda_1 + \lambda_2) \right] \pi(0, 0)
        \end{split} \nonumber \\
        \intertext{which simplifies to}
        \begin{split}
            & \left[ (\lambda_1 + \lambda_2 + \mu) - \mu y^{-1} - \lambda_2 y \right] P_y(y)
            + \left[ \gamma_2 y + \theta_2(y-1) \right] \frac{d}{dy} P_y(y) \\
            & \quad = \quad \left[ \mu + \gamma_1 y + \theta_1 \right] P_y(y | n_1=1)
            + \left[ \mu - \mu y^{-1} \right] \pi(0, 0)
        \end{split} \nonumber \\
        \intertext{and, multiplying through by $xy$:}
        \begin{split}
            & \left[ (\lambda_1 + \lambda_2 + \mu)xy - \mu x - \lambda_2 x y^2 \right] P_y(y)
            + xy \left[ \gamma_2 y + \theta_2(y-1) \right] \frac{d}{dy} P_y(y) \\
            & \quad = \quad \left[ \mu + \gamma_1 y + \theta_1 \right] xy \, P_y(y | n_1=1)
            - \mu x(1-y) \, \pi(0,0)
        \end{split} \label{eq:gen:transform_y_boundary}
    \end{align}
    
    Substituting Equation~\eqref{eq:gen:transform_x_boundary} into Equation~\eqref{eq:gen:transform_interior} eliminates all $P_x(x)$ terms:
    \begin{align}
        \begin{split}
            & \left[ (\lambda_1+\lambda_2+\mu)xy - \mu y - \lambda_1 x^2 y - \lambda_2 x y^2 \right] P(x,y) \\
            & + xy \left[ \gamma_1(x-y) + \theta_1(x-1) \right] \frac{\partial}{\partial x} P(x,y) \\
            & + xy \left[ \gamma_2(y-x) + \theta_2(y-1) \right] \frac{\partial}{\partial y} P(x,y) \\
            & \quad = \quad \left[ (\lambda_1 + \lambda_2 + \mu)xy - \mu y - \lambda_2 x y^2 \right] P_y(y) \\
            & \quad \quad + xy \left[ \gamma_2 y + \theta_2(y-1) \right] \frac{d}{dy} P_y(y) \\
            & \quad \quad - xy \left[ \mu + \gamma_1 y + \theta_1 \right] P_y(y | n_1=1)
        \end{split} \label{eq:gen:plugged_interior_x}
    \end{align}
    
    Isolating $[\mu+\gamma_1 y+\theta_1]\,xy\,P_y(y|n_1=1)$ from Equation~\eqref{eq:gen:transform_y_boundary} and substituting into Equation~\eqref{eq:gen:plugged_interior_x}:
    \begin{align*}
        \begin{split}
            & \left[ (\lambda_1+\lambda_2+\mu)xy - \mu y - \lambda_1 x^2 y - \lambda_2 x y^2 \right] P(x,y) \\
            & + xy \left[ \gamma_1(x-y) + \theta_1(x-1) \right] \frac{\partial}{\partial x} P(x,y) \\
            & + xy \left[ \gamma_2(y-x) + \theta_2(y-1) \right] \frac{\partial}{\partial y} P(x,y) \\
            & \quad = \quad \left[ (\lambda_1 + \lambda_2 + \mu)xy - \mu y - \lambda_2 x y^2 \right] P_y(y) \\
            & \quad \quad + xy \left[ \gamma_2 y + \theta_2(y-1) \right] \frac{d}{dy} P_y(y) \\
            & \quad \quad - \left[ (\lambda_1 + \lambda_2 + \mu)xy - \mu x - \lambda_2 x y^2 \right] P_y(y) \\
            & \quad \quad - xy \left[ \gamma_2 y + \theta_2(y-1) \right] \frac{d}{dy} P_y(y)
            \;-\; \mu x(1-y) \, \pi(0,0)
        \end{split} \nonumber \\
        \intertext{The two $\frac{d}{dy}P_y(y)$ terms cancel and the two $P_y(y)$ coefficients combine into $\mu(x-y)$, leaving}
        \begin{split}
            & \left[ (\lambda_1+\lambda_2+\mu)xy - \mu y - \lambda_1 x^2 y - \lambda_2 x y^2 \right] P(x,y) \\
            & + xy \left[ \gamma_1(x-y) + \theta_1(x-1) \right] \frac{\partial}{\partial x} P(x,y) \\
            & + xy \left[ \gamma_2(y-x) + \theta_2(y-1) \right] \frac{\partial}{\partial y} P(x,y) \\
            & \quad = \quad \mu(x-y) P_y(y) - \mu x(1-y) \, \pi(0,0)
        \end{split}
    \end{align*}
    This is Equation~\eqref{eq:gen:fundamental_equation}, which completes the proof.
\end{proof}

The fundamental equation also delivers a diagonal identity in the no-abandonment case, complementing Corollary~\ref{cor:gen:no_abandon}.

\begin{corollary}[Diagonal identity without abandonment]\label{cor:gen:diagonal}
    Suppose $\theta_1=\theta_2=0$ (Models~$A$, $B$, $B_2$ and $\BH$) and $\rho<1$. Setting $x=y=z$ in the fundamental equation~\eqref{eq:gen:fundamental_equation} annihilates both derivative terms (each carries the factor $x-y$), leaving the algebraic relation
    \begin{equation*}
        -(1-z)\bigl(\mu-\lambda z\bigr)P(z,z)=-\mu(1-z)\,\pi(0,0),
    \end{equation*}
    whence the diagonal of the PGF and the normalisation read
    \begin{equation}
        P(z,z)=\frac{\pi(0,0)}{1-\rho z}, \qquad P(1,1)=1-\pi_0=\rho.
        \label{eq:gen:Pzz_diag}
    \end{equation}
    Since $\pi(0,0)=\rho(1-\rho)$ by Corollary~\ref{cor:gen:no_abandon}, the diagonal is, up to the busy probability $\rho$, the PGF of a geometric distribution: conditionally on the server being busy, the total number of waiting customers is geometric, $\mathbb{P}(N=n \mid \mathcal{B}=1)=(1-\rho)\rho^{n}$ for $n\ge0$.
\end{corollary}


\subsubsection{Equation for PPGF in state space \texorpdfstring{$\widetilde{S}$}{S~}}

The equation of Lemma~\ref{lem:gen:fundamental_equation} transforms both coordinates at once. An alternative device, notably used by Cohen~\cite{cohen1956delay}, is the \textbf{Partial Probability Generating Function (PPGF)}: transforming only a subset of the random variables while keeping the discrete indices of the rest. For a priority queue this is particularly advantageous: it lets us exploit the well-known properties of the priority class (often a standard $M/M/1$ system) and concentrate the analytical effort on the non-trivial PGF of the non-priority class. Cohen developed this framework for multi-server systems and for the state space $S$; we apply it here to the single-server case, in a modern formulation consistent with the variables defined above, and ---in this subsection--- on the alternative state space $\widetilde{S}$.

On $S$, the PPGF with respect to $N_2$, jointly with the event that the first queue holds exactly $n_1$ waiting customers, is
\begin{equation}
    \widetilde{P}(n_1, y) = \mathbb{E} \left[y^{N_2} \mathbf{1}_{\{ N_1=n_1\}} \right] = \sum_{n_2=0}^{\infty} \pi(n_1,n_2) y^{n_2},
    \label{eq:gen:ppgf_n1}
\end{equation}
and these partial transforms recover the joint PGF through $P(x,y) = \sum_{n_1=0}^{\infty} \widetilde{P}(n_1, y)\, x^{n_1}$. Despite the conditional look of the indicator, Equation~\eqref{eq:gen:ppgf_n1} is a joint (unnormalised) transform; this version is used in the Model-$A$ analysis (\S\ref{sec:model_a}). On $\widetilde{S}$, the analogous object transforms $n_2$ and keeps the total count $n$ discrete:
\begin{equation}
    \widetilde{P}(y,n)=\sum_{n_2=0}^{n}\tpi(n_2,n)\,y^{n_2},
    \label{eq:gen:ppgf_Stilde}
\end{equation}
a \emph{polynomial} of degree $n$ in $y$, since $n_2\le n$ on $\widetilde{S}$; this finiteness justifies term-by-term differentiation and boundedness arguments without convergence caveats.

\begin{lemma}[Fundamental equation for the PPGF on $\widetilde{S}$]\label{lem:gen:PPGF_dynamics}
    Consider the same system on the state space $\widetilde{S}$ of
    Equation~\eqref{eq:gen:state_space_Stilde}, where $n$ is the total number of customers
    \emph{waiting} in the queues and $n_2$ the number of waiting class-$2$ customers.
    Assume positive recurrence, so that the stationary distribution
    $\{\tpi(n_2,n)\}$ exists, and let $\widetilde{P}(y,n)$ be the partial PGF of
    Equation~\eqref{eq:gen:ppgf_Stilde}. Then, for every $n\ge 1$, the family
    $\{\widetilde{P}(y,n)\}_{n\ge 0}$ satisfies the differential--difference equation
    \begin{equation*}
    \begin{split}
        & \left[ -(\gamma_2 + \gamma_1 y)(1-y) + (\theta_2 - \theta_1) y \right] \frac{d}{dy} \widetilde{P}(y,n)
          + \left[ \lambda_1 + \lambda_2 + \mu + \gamma_1 n(1-y) + \theta_1 n \right] \widetilde{P}(y,n) \\
        & \quad = \quad (\lambda_1 + \lambda_2 y) \widetilde{P}(y,n-1)
          + \left[ \mu + \theta_1 (n+1) \right] \widetilde{P}(y,n+1) \\
        & \quad \quad + (\theta_2 - \theta_1 y) \frac{d}{dy} \widetilde{P}(y,n+1)
          + \mu y^{n} (1-y)\, \tpi(n+1,n+1).
    \end{split}
    \end{equation*}
\end{lemma}

\begin{proof}
    Multiply the interior equation~\eqref{eq:gen:tilde_balance_interior} by $y^{n_2}$ and sum over $1\le n_2\le n-1$:
    \begin{align*}
        \begin{split}
            & (\lambda_1+\lambda_2+\mu) \sum_{n_2=1}^{n-1} \tpi(n_2,n) y^{n_2} \\
            & + \gamma_1 \sum_{n_2=1}^{n-1} (n-n_2) \tpi(n_2,n) y^{n_2} \\
            & + \gamma_2 \sum_{n_2=1}^{n-1} n_2 \tpi(n_2,n) y^{n_2} \\
            & + \theta_1 \sum_{n_2=1}^{n-1} (n-n_2) \tpi(n_2,n) y^{n_2} \\
            & + \theta_2 \sum_{n_2=1}^{n-1} n_2 \tpi(n_2,n) y^{n_2} \\
            & \quad = \quad \lambda_1 \sum_{n_2=1}^{n-1} \tpi(n_2, n-1) y^{n_2} \\
            & \quad \quad + \lambda_2 \sum_{n_2=1}^{n-1} \tpi(n_2-1, n-1) y^{n_2} \\
            & \quad \quad + \mu \sum_{n_2=1}^{n-1} \tpi(n_2, n+1) y^{n_2} \\
            & \quad \quad + \gamma_1 \sum_{n_2=1}^{n-1} (n-n_2+1) \tpi(n_2-1, n) y^{n_2} \\
            & \quad \quad + \gamma_2 \sum_{n_2=1}^{n-1} (n_2+1) \tpi(n_2+1, n) y^{n_2} \\
            & \quad \quad + \theta_1 \sum_{n_2=1}^{n-1} (n+1-n_2) \tpi(n_2, n+1) y^{n_2} \\
            & \quad \quad + \theta_2 \sum_{n_2=1}^{n-1} (n_2+1) \tpi(n_2+1, n+1) y^{n_2}
        \end{split}
    \end{align*}
    
    Adding the diagonal equation~\eqref{eq:gen:tilde_balance_diagonal} multiplied by $y^n$ extends most summation limits to $n$; the $\lambda_1$- and $\gamma_2$-terms on the RHS remain restricted to $n_2\le n-1$, and an extra term $\mu\tpi(n+1,n+1)y^n$ remains:
    \begin{align*}
        \begin{split}
            & (\lambda_1+\lambda_2+\mu) \sum_{n_2=1}^{n} \tpi(n_2,n) y^{n_2} \\
            & + \gamma_1 \sum_{n_2=1}^{n} (n-n_2) \tpi(n_2,n) y^{n_2} \\
            & + \gamma_2 \sum_{n_2=1}^{n} n_2 \tpi(n_2,n) y^{n_2} \\
            & + \theta_1 \sum_{n_2=1}^{n} (n-n_2) \tpi(n_2,n) y^{n_2} \\
            & + \theta_2 \sum_{n_2=1}^{n} n_2 \tpi(n_2,n) y^{n_2} \\
            & \quad = \quad \lambda_1 \sum_{n_2=1}^{n-1} \tpi(n_2,n-1) y^{n_2} \\
            & \quad \quad + \lambda_2 \sum_{n_2=1}^{n} \tpi(n_2-1,n-1) y^{n_2} \\
            & \quad \quad + \mu \sum_{n_2=1}^{n} \tpi(n_2,n+1) y^{n_2} + \mu \tpi(n+1,n+1) y^n \\
            & \quad \quad + \gamma_1 \sum_{n_2=1}^{n} (n-n_2+1) \tpi(n_2-1,n) y^{n_2} \\
            & \quad \quad + \gamma_2 \sum_{n_2=1}^{n-1} (n_2+1) \tpi(n_2+1,n) y^{n_2} \\
            & \quad \quad + \theta_1 \sum_{n_2=1}^{n} (n+1-n_2) \tpi(n_2,n+1) y^{n_2} \\
            & \quad \quad + \theta_2 \sum_{n_2=1}^{n} (n_2+1) \tpi(n_2+1,n+1) y^{n_2}.
        \end{split}
    \end{align*}

    As in the proof of Lemma~\ref{lem:gen:fundamental_equation}, we write $LHS_i$ and $RHS_j$ for the $i$-th and $j$-th summands on the left- and right-hand sides of this display, in the order written (the two $\mu$-terms together forming $RHS_3$). Evaluating each summand:
    \begin{align*}
        \begin{split}
            & LHS_1 \\
            & \quad = \quad (\lambda_1 + \lambda_2 + \mu) \sum_{n_2=1}^{n} \tpi(n_2,n) y^{n_2} \\
            & \quad = \quad (\lambda_1 + \lambda_2 + \mu) \left[ \sum_{n_2=0}^{n} \tpi(n_2,n) y^{n_2} - \tpi(0,n) y^0 \right] \\
            & \quad = \quad (\lambda_1 + \lambda_2 + \mu) \left[ \widetilde{P}(y,n) - \tpi(0,n) \right]
        \end{split} \\
        \begin{split}
            & LHS_2 \\
            & \quad = \quad \gamma_1 \sum_{n_2=1}^{n} (n-n_2) \tpi(n_2,n) y^{n_2} \\
            & \quad = \quad \gamma_1 \left[ n \sum_{n_2=1}^{n} \tpi(n_2,n) y^{n_2} - \sum_{n_2=1}^{n} n_2 \tpi(n_2,n) y^{n_2} \right] \\
            & \quad = \quad \gamma_1 \left[ n \left( \sum_{n_2=0}^{n} \tpi(n_2,n) y^{n_2} - \tpi(0,n) y^0 \right) - y \sum_{n_2=1}^{n} n_2 \tpi(n_2,n) y^{n_2-1} \right] \\
            & \quad = \quad \gamma_1 \left[ n \big( \widetilde{P}(y,n) - \tpi(0,n) \big) - y \frac{d}{dy} \widetilde{P}(y,n) \right]
        \end{split} \\
        \begin{split}
            & LHS_3 \\
            & \quad = \quad \gamma_2 \sum_{n_2=1}^{n} n_2 \tpi(n_2,n) y^{n_2} \\
            & \quad = \quad \gamma_2 \left[ y \sum_{n_2=1}^{n} n_2 \tpi(n_2,n) y^{n_2-1} \right] \\
            & \quad = \quad \gamma_2 y \frac{d}{dy} \widetilde{P}(y,n)
        \end{split} \\
        \begin{split}
            & LHS_4 \\
            & \quad = \quad \theta_1 \sum_{n_2=1}^{n} (n-n_2) \tpi(n_2,n) y^{n_2} \\
            & \quad = \quad \theta_1 \left[ n \big( \widetilde{P}(y,n) - \tpi(0,n) \big) - y \frac{d}{dy} \widetilde{P}(y,n) \right]
        \end{split} \\
        \begin{split}
            & LHS_5 \\
            & \quad = \quad \theta_2 \sum_{n_2=1}^{n} n_2 \tpi(n_2,n) y^{n_2} \\
            & \quad = \quad \theta_2 y \frac{d}{dy} \widetilde{P}(y,n)
        \end{split} \\
        \begin{split}
            & RHS_1 \\
            & \quad = \quad \lambda_1 \sum_{n_2=1}^{n-1} \tpi(n_2,n-1) y^{n_2} \\
            & \quad = \quad \lambda_1 \left[ \sum_{n_2=1}^{n} \tpi(n_2,n-1) y^{n_2} - \tpi(n,n-1) y^n \right] \\
            & \quad = \quad \lambda_1 \left[ \sum_{n_2=0}^{n} \tpi(n_2,n-1) y^{n_2} - \tpi(n,n-1) y^n - \tpi(0,n-1) y^0 \right] \\
            & \quad = \quad \lambda_1 \left[ \widetilde{P}(y,n-1) - \tpi(n,n-1) y^n - \tpi(0,n-1) \right]
        \end{split} \\
        \begin{split}
            & RHS_2 \\
            & \quad = \quad \lambda_2 \sum_{n_2=1}^{n} \tpi(n_2-1,n-1) y^{n_2} \\
            & \quad = \quad \lambda_2 \left[ y \sum_{n_2=1}^{n} \tpi(n_2-1,n-1) y^{n_2-1} \right] \\
            & \quad = \quad \lambda_2 \left[ y \sum_{n_2=0}^{n-1} \tpi(n_2,n-1) y^{n_2} \right] \\
            & \quad = \quad \lambda_2 \left[ y \left( \sum_{n_2=0}^{n} \tpi(n_2,n-1) y^{n_2} - \tpi(n,n-1) y^n \right) \right] \\
            & \quad = \quad \lambda_2 y \left[ \widetilde{P}(y,n-1) - \tpi(n,n-1) y^n \right]
        \end{split} \\
        \begin{split}
            & RHS_3 \\
            & \quad = \quad \mu \tpi(n+1,n+1) y^n + \mu \sum_{n_2=1}^{n} \tpi(n_2,n+1) y^{n_2} \\
            & \quad = \quad \mu \left[ \tpi(n+1,n+1) y^n + \sum_{n_2=1}^{n} \tpi(n_2,n+1) y^{n_2} \right] \\
            & \quad = \quad \mu \left[ \tpi(n+1,n+1) y^n + \sum_{n_2=0}^{n+1} \tpi(n_2,n+1) y^{n_2} - \tpi(0,n+1) y^0 - \tpi(n+1,n+1) y^{n+1} \right] \\
            & \quad = \quad \mu \left[ \widetilde{P}(y,n+1) - \tpi(0,n+1) + \tpi(n+1,n+1) y^n (1-y) \right]
        \end{split} \\
        \begin{split}
            & RHS_4 \\
            & \quad = \quad \gamma_1 \sum_{n_2=1}^{n} (n-n_2+1) \tpi(n_2-1,n) y^{n_2} \\
            & \quad = \quad \gamma_1 \left[ y \sum_{n_2=1}^{n} (n - (n_2-1)) \tpi(n_2-1,n) y^{n_2-1} \right] \\
            & \quad = \quad \gamma_1 \left[ y \sum_{k=0}^{n-1} (n-k) \tpi(k,n) y^k \right] \\
            & \quad = \quad \gamma_1 \left[ y \sum_{k=0}^{n} (n-k) \tpi(k,n) y^k \right] \quad \text{(since the $k=n$ term is 0)} \\
            & \quad = \quad \gamma_1 y \left[ n \sum_{k=0}^{n} \tpi(k,n) y^k - y \sum_{k=1}^{n} k \tpi(k,n) y^{k-1} \right] \\
            & \quad = \quad \gamma_1 y \left[ n \widetilde{P}(y,n) - y \frac{d}{dy} \widetilde{P}(y,n) \right]
        \end{split} \\
        \begin{split}
            & RHS_5 \\
            & \quad = \quad \gamma_2 \sum_{n_2=1}^{n-1} (n_2+1) \tpi(n_2+1,n) y^{n_2} \\
            & \quad = \quad \gamma_2 \sum_{k=2}^{n} k \tpi(k,n) y^{k-1} \quad \text{(letting $k = n_2+1$)} \\
            & \quad = \quad \gamma_2 \left[ \sum_{k=1}^{n} k \tpi(k,n) y^{k-1} - 1 \cdot \tpi(1,n) y^0 \right] \\
            & \quad = \quad \gamma_2 \left[ \frac{d}{dy} \widetilde{P}(y,n) - \tpi(1,n) \right]
        \end{split} \\
        \begin{split}
            & RHS_6 \\
            & \quad = \quad \theta_1 \sum_{n_2=1}^{n} (n+1-n_2) \tpi(n_2,n+1) y^{n_2} \\
            & \quad = \quad \theta_1 \left[ (n+1) \sum_{n_2=1}^{n} \tpi(n_2,n+1) y^{n_2} - \sum_{n_2=1}^{n} n_2 \tpi(n_2,n+1) y^{n_2} \right] \\
            & \quad = \quad \theta_1 \left[ (n+1) \left( \sum_{n_2=0}^{n+1} \tpi(n_2,n+1) y^{n_2} - \tpi(0,n+1) - \tpi(n+1,n+1) y^{n+1} \right) \right] \\
            & \quad \quad - \theta_1 \left[ y \left( \sum_{n_2=0}^{n+1} n_2 \tpi(n_2,n+1) y^{n_2-1} - (n+1) \tpi(n+1,n+1) y^n \right) \right] \\
            & \quad = \quad \theta_1 \left[ (n+1) \big( \widetilde{P}(y,n+1) - \tpi(0,n+1) \big) - (n+1) \tpi(n+1,n+1) y^{n+1} \right] \\
            & \quad \quad - \theta_1 \left[ y \frac{d}{dy} \widetilde{P}(y,n+1) - (n+1) \tpi(n+1,n+1) y^{n+1} \right] \\
            & \quad = \quad \theta_1 \left[ (n+1) \big( \widetilde{P}(y,n+1) - \tpi(0,n+1) \big) - y \frac{d}{dy} \widetilde{P}(y,n+1) \right]
        \end{split} \\
        \begin{split}
            & RHS_7 \\
            & \quad = \quad \theta_2 \sum_{n_2=1}^{n} (n_2+1) \tpi(n_2+1,n+1) y^{n_2} \\
            & \quad = \quad \theta_2 \sum_{k=2}^{n+1} k \tpi(k,n+1) y^{k-1} \quad \text{(letting $k = n_2+1$)} \\
            & \quad = \quad \theta_2 \left[ \sum_{k=1}^{n+1} k \tpi(k,n+1) y^{k-1} - 1 \cdot \tpi(1,n+1) y^0 \right] \\
            & \quad = \quad \theta_2 \left[ \frac{d}{dy} \widetilde{P}(y,n+1) - \tpi(1,n+1) \right]
        \end{split}
    \end{align*}
    
    Combining all the terms above:
    \begin{align*}
        \begin{split}
            & (\lambda_1 + \lambda_2 + \mu) \left[ \widetilde{P}(y,n) - \tpi(0,n) \right] \\
            & + \gamma_1 \left[ n \big( \widetilde{P}(y,n) - \tpi(0,n) \big) - y \frac{d}{dy} \widetilde{P}(y,n) \right] \\
            & + \gamma_2 y \frac{d}{dy} \widetilde{P}(y,n) \\
            & + \theta_1 \left[ n \big( \widetilde{P}(y,n) - \tpi(0,n) \big) - y \frac{d}{dy} \widetilde{P}(y,n) \right] \\
            & + \theta_2 y \frac{d}{dy} \widetilde{P}(y,n) \\
            & \quad = \quad \lambda_1 \left[ \widetilde{P}(y,n-1) - \tpi(n,n-1) y^n - \tpi(0,n-1) \right] \\
            & \quad \quad + \lambda_2 y \left[ \widetilde{P}(y,n-1) - \tpi(n,n-1) y^n \right] \\
            & \quad \quad + \mu \left[ \widetilde{P}(y,n+1) - \tpi(0,n+1) + \tpi(n+1,n+1) y^n (1-y) \right] \\
            & \quad \quad + \gamma_1 y \left[ n \widetilde{P}(y,n) - y \frac{d}{dy} \widetilde{P}(y,n) \right] \\
            & \quad \quad + \gamma_2 \left[ \frac{d}{dy} \widetilde{P}(y,n) - \tpi(1,n) \right] \\
            & \quad \quad + \theta_1 \left[ (n+1) \big( \widetilde{P}(y,n+1) - \tpi(0,n+1) \big) - y \frac{d}{dy} \widetilde{P}(y,n+1) \right] \\
            & \quad \quad + \theta_2 \left[ \frac{d}{dy} \widetilde{P}(y,n+1) - \tpi(1,n+1) \right]
        \end{split}
    \end{align*}
    Grouping the common terms yields:
    \begin{align*}
        \begin{split}
            & \left[ -(\gamma_2 + \gamma_1 y)(1-y) + (\theta_2 - \theta_1) y \right] \frac{d}{dy} \widetilde{P}(y,n) + \left[ \lambda_1 + \lambda_2 + \mu + \gamma_1 n(1-y) + \theta_1 n \right] \widetilde{P}(y,n) \\
            & \quad = \quad (\lambda_1 + \lambda_2 y) \widetilde{P}(y,n-1) + \left[ \mu + \theta_1 (n+1) \right] \widetilde{P}(y,n+1) + (\theta_2 - \theta_1 y) \frac{d}{dy} \widetilde{P}(y,n+1) \\
            & \quad \quad + \left\{ \left[ \lambda_1 + \lambda_2 + \mu + (\gamma_1 + \theta_1) n \right] \tpi(0,n) - \lambda_1 \tpi(0,n-1) - \left[ \mu + \theta_1 (n+1) \right] \tpi(0,n+1) \right. \\
            & \quad \quad \quad \left. - \gamma_2 \tpi(1,n) - \theta_2 \tpi(1,n+1) \right\} \\
            & \quad \quad - y^n (\lambda_1 + \lambda_2 y) \tpi(n,n-1) + \mu y^n (1-y) \tpi(n+1,n+1)
        \end{split}
    \end{align*}
    The $y$-boundary equation~\eqref{eq:gen:tilde_balance_y_boundary} makes the bracketed terms vanish. Moreover, $\tpi(n,n-1)=0$, because the state $(n,n-1)$ lies outside $\widetilde{S}$. Hence
    \begin{align*}
        \begin{split}
            & \left[ -(\gamma_2 + \gamma_1 y)(1-y) + (\theta_2 - \theta_1) y \right] \frac{d}{dy} \widetilde{P}(y,n) + \left[ \lambda_1 + \lambda_2 + \mu + \gamma_1 n(1-y) + \theta_1 n \right] \widetilde{P}(y,n) \\
            & \quad = \quad (\lambda_1 + \lambda_2 y) \widetilde{P}(y,n-1) + \left[ \mu + \theta_1 (n+1) \right] \widetilde{P}(y,n+1) + (\theta_2 - \theta_1 y) \frac{d}{dy} \widetilde{P}(y,n+1) \\
            & \quad \quad + \mu y^n (1-y) \tpi(n+1,n+1)
        \end{split}
    \end{align*}
    This is the asserted differential--difference equation, which completes the proof.
\end{proof}

\begin{remark}[The $\widetilde{S}$ recursion is a failed simplification]\label{rem:gen:Stilde_failed}
The partial-PGF recurrence of Lemma~\ref{lem:gen:PPGF_dynamics} was introduced with the
aim that taking a generating function over the class-$2$ index only, while carrying the total
count $n$ as a discrete parameter, would \emph{simplify} the analysis~---~collapsing the
two-dimensional balance on $S$ to a one-dimensional recurrence in $n$. The opposite turns out
to be true. The inhomogeneous term $\mu y^{n}(1-y)\,\tpi(n+1,n+1)$ couples each level $n$ of
the recurrence to the diagonal probability $\tpi(n+1,n+1)=\pi(0,n+1)$, and these diagonal
probabilities are themselves unknown. The recurrence therefore does not close: solving it
would require exactly the boundary probabilities it was intended to produce. Only for the
baseline Model-$A$ does this coupling term decay fast enough to be neglected, and even there
dropping it yields only an approximate solution, accurate in a restricted parameter region
(Approximation~\ref{apx:A:PPGF_Stilde}). For every model with an active mechanism the
$\widetilde{S}$ route is strictly harder than the $S$ analysis it was meant to replace. We
therefore record it as a deliberate negative result: the $\widetilde{S}$ coordinatisation is
retained only for Model-$A$, and the $S$ formulation is used throughout the remainder of this
thesis.
\end{remark}


